% !TeX root = ../../Book4.tex
% 纲要标题：## 第4章　周期、Abel 定理与 Jacobi 簇
% 纲要位置：几何学计划纲要.md，第 3588 行。
% 正文状态：初稿；尚未正式审核。

% 纲要内容开始
% 
% > 本章目标与证明范围：由具体周期积分构造 Jacobi 簇，再证明主除子的判据和反演。椭圆函数是完整低亏格模型，不仅作为一条跨学科联系。
% 
% 纲要内容结束

\chapter{周期、Abel 定理与 Jacobi 簇}
\label{g4:ch:04}
\BookChapterNumber{g4:ch:04}

\begin{chapterabstract}
	亚纯函数的零极点总次数为零，但这个条件在正亏格曲面上并不充分。本章把全纯微分沿闭路的积分组织成周期矩阵，由其对称性与正性构造 Jacobi 簇。第三类微分将给定除子转化为对数微分的候选；它的整数周期条件给出 Abel 定理。Riemann--Roch 与映射度继而证明 Jacobi 反演，并把反演的全部解解释为完全线性系。theta 级数构造标准主极化；有效除子像与 theta 除子的识别，以及由主极化恢复曲面的 Torelli 定理，采用明确的文献陈述。最后从格级数完整建立椭圆函数、光滑三次模型、弦切群律和 \(j\) 分类，并计算一个周期趋于无穷远时的结点退化。
\end{chapterabstract}

\begin{learninggoals}
	\begin{itemize}
		\item 从切割多边形推导周期双线性公式，证明归一化周期矩阵的对称性和虚部正定性。
		\item 构造 Jacobi 环面、Abel--Jacobi 映射与对称积上的积分映射，并计算重合点处的微分。
		\item 构造归一化第三类微分，核对带极点互反公式的符号，以指数积分证明主除子的完整判据。
		\item 用 Riemann--Roch 和映射度证明反演满射性，计算特殊除子对应的纤维与微分秩。
		\item 验证 theta 的收敛、变换律、截面维数和极化形式，辨认 theta 除子识别及 Torelli 重构所用的外部定理。
		\item 从 Weierstrass 级数推导三次方程和加法公式，用格变换、\(j\) 不变量及退化极限比较椭圆曲线。
	\end{itemize}
\end{learninggoals}

在球面上，两个不同点 \(p,q\) 总能成为某个亚纯函数唯一的单零点和单极点；在有限坐标中，函数 \((z-p)/(z-q)\) 就完成了这件事。在正亏格曲面上，同样的要求却不可能成立：这样的函数到球面的次数为 \(1\)，会把曲面本身变成球面。把问题改为寻找除子 \(n[p]-n[q]\) 的函数，答案开始依赖两点之间的积分。在环面 \(\mathbb C/\Lambda\) 上，充要条件将是 \(n(p-q)=0\)；某些点对可以实现，另一些不能。次数为零之后还剩下的条件，正由周期记录。

\ProseParagraphBreak

Abel 定理\footnote{Niels Henrik Abel（1802--1829），挪威数学家，研究椭圆函数与代数函数的积分。}把这一现象推广到任意紧黎曼曲面。先对全纯微分沿同调辛基积分，得到一个格；再把点之间的积分模去这个格，便得到除子的整体不变量。Jacobi 簇\footnote{Carl Gustav Jacob Jacobi（1804--1851），德国数学家，研究椭圆函数、theta 函数与动力学。}容纳全部这些积分值，反演定理说明每个值都能由 \(g\) 个点实现。实现的方式未必唯一，而所有实现组成的射影空间恰好就是第三章的完全线性系。

\ProseParagraphBreak

亏格一时只有一条独立的全纯微分，曲面经积分就成为自己的 Jacobi 环面。这时可以逐项求和，把抽象存在的函数直接写出来：Weierstrass 函数\footnote{Karl Theodor Wilhelm Weierstraß（1815--1897），德国数学家，研究分析的严格基础与复函数论。}在每个格点放置一个二阶极点，它的导数提供三阶极点。消去 Laurent 主部后，二者满足光滑三次方程；复数加法则变为三次曲线上的弦切作图。这个例子也让周期矩阵、除子、积分和代数方程在同一个对象上相遇。

\ProseParagraphBreak

本章直接使用第三章的主部判据、Riemann--Roch 及典范微分维数，局部分析仍是第一章的 Cauchy 积分与留数。拓扑计算所需的曲面切割、相交数和映射度在调用处恢复其含义与约定。\secref{g4:sec:0401}与\secref{g4:sec:0403}分别给出无极点及有极点的双线性计算；\secref{g4:sec:0404}证明反演与纤维公式，再把 Riemann 的 theta 定理和经典 Torelli 定理作为有精确出处的外部结果。有关 Jacobi 簇与 Abel 定理的阅读导航见\parencite[第20--21节]{Forster1981}；周期、theta 及曲线 Jacobi 簇的进一步讨论见\parencite[第2章第6--7节]{GriffithsHarris1994}。

% !TeX root = ../../Book4.tex
% 纲要标题：### 4.1 微分的周期
% 纲要位置：几何学计划纲要.md，第 3592 行。
% 正文状态：初稿；尚未正式审核。

% 纲要内容开始
% 
% * 同调辛基与周期矩阵
% * Riemann 双线性关系
% * 归一化全纯微分
% * 周期格的离散性与正性
% 
% 纲要内容结束

\section{微分的周期}
\label{g4:sec:0401}

全纯微分局部可以积分成全纯函数，但沿闭路积分未必为零。例如在复环面上，\(\mathrm{d}z\) 是全局微分，沿格向量对应的闭路积分正好得到该格向量。要把这种现象推广到任意紧黎曼曲面，需要同时记录曲面的闭路和微分的积分。第三章已证明全纯微分空间的复维数等于拓扑亏格；本节进一步说明，这些积分组成的矩阵满足对称性与正性，并由此产生一个离散的周期格。

\subsection{同调辛基与周期矩阵}
\label{g4:subsec:0401:01}

在有向曲面上，两条横截的有向闭曲线在每个交点给出一个符号：若先取第一条、再取第二条的正向切向量得到曲面的正定向，记为 \(+1\)，否则记为 \(-1\)。符号之和称为它们的相交数。第二卷定理16.2.8及推论16.2.9证明，对于闭有向光滑曲面，这个数只依赖整数同调类，且给出反对称双线性配对。分片光滑曲线可先在有限个折点附近光滑化；小扰动不改变其同调类。因此以下把相交数写成 \(a\cdot b\)，并固定顺序为先 \(a\) 后 \(b\)。

\begin{proposition}
	\label{g4:prop:0401:symplectic-basis}
	设 \(X\) 是亏格为 \(g\) 的紧连通黎曼曲面。则 \(H_1(X;\mathbb Z)\) 是秩为 \(2g\) 的自由阿贝尔群，并存在一组基
	\[
		a_1,\ldots,a_g,b_1,\ldots,b_g
	\]
	满足
	\begin{equation}\label{g4:eq:0401:symplectic-pairing}
		a_i\cdot a_j=b_i\cdot b_j=0,
		\qquad a_i\cdot b_j=\delta_{ij}.
	\end{equation}
	当 \(g\ge1\) 时，可取这些类的分片光滑闭路代表，使沿它们切开 \(X\) 后得到一张闭多边形 \(P\)，其正向边界字为
	\begin{equation}\label{g4:eq:0401:polygon-word}
		a_1b_1a_1^{-1}b_1^{-1}\cdots
		a_gb_ga_g^{-1}b_g^{-1}.
	\end{equation}
	多边形内部映到曲面的切口补集，各配对边按相同字母识别，所有顶点映到同一点。切口还可避开预先指定的有限个点。
\end{proposition}
\begin{proof}
	黎曼曲面的坐标变换保持实定向，故 \(X\) 是闭有向光滑曲面。第一卷命题6.4.5的可定向标准字构造给出式\eqref{g4:eq:0401:polygon-word}的拓扑切割模型；亏格的唯一性来自该卷定理6.5.4。这个模型有一个顶点、\(2g\) 条边和一个二维胞腔。胞腔边界 \(\partial_1\) 为零；\(\partial_2\) 的每个系数也为零，因为每条边在面边界中正、反各出现一次。第二卷定理3.1.6将胞腔同调自然识别为奇异同调，因而
	\[
		H_1(X;\mathbb Z)=\mathbb Z a_1\oplus\cdots\oplus\mathbb Z a_g
		\oplus\mathbb Z b_1\oplus\cdots\oplus\mathbb Z b_g.
	\]
	在每个手柄中取经线与纬线，安排先经线、后纬线的切向量为正向，则它们相交一次且符号为正；不同手柄中的核心曲线不交。同一条曲线可向旁边平移成与自己不交的代表。它们与上述边类对应，故得到式\eqref{g4:eq:0401:symplectic-pairing}。\(g=0\) 时用球面的一个零胞腔和一个二胞腔即可得到 \(H_1=0\)。

	这里的拓扑切口还须变成可积分的曲线。第一卷引理6.2.6及其后的有限图邻域构造，为切口图提供互不相交的顶点小盘和截短边的窄带。沿每条窄带中的核心弧取有限条光滑坐标小盘链，在相邻小盘的交集中选连接点，以坐标线段相接。各段在黎曼曲面上实解析，稍移连接点以避开共同子弧后，紧的有限段只有有限次交叉；删去回路，便得到仍在该带内、端点相同的简单分片光滑弧。顶点附近先在光滑坐标中画出具有原循环次序的有限条短射线，再在各扇区内连接到这些边弧。不同边使用不交的窄带，故没有引入新的相交。

	这个替换保持切割类型。其理由与第一卷引理6.2.12的证明相同：在顶点盘的各扇区及边带内，适当简单弧的延拓定理把旧弧送到新弧，并固定支撑盘的边界；各盘内的同胚再延拓为环境同痕。因而补集仍是一张盘，边字不变。可在顶点处令相邻射线不共线，边的其余折点再作局部圆滑，使切开的区域能分成有限个光滑带角块。若有有限个指定点，先在各点的小坐标盘内把经过该点的有限条弧绕开，再实施同样的替换即可。
\end{proof}

满足式\eqref{g4:eq:0401:symplectic-pairing}的整数基称为同调的\term[tongdiaoxinji]{辛基}[symplectic basis]。一个辛基首先是同调类的基；为了进行边界积分，我们另外选择了具有式\eqref{g4:eq:0401:polygon-word}的切口代表。之后的积分结论只依赖这些同调类。

\ProseParagraphBreak

若 \(\alpha\) 是 \(X\) 上光滑的复值闭一形式，沿闭路的积分只依赖整数同调类。这里所需的桥梁是第二卷定理7.3.5：在无边界光滑流形上，光滑奇异链计算与连续奇异链相同的同调；相同同调类的光滑圈之差因而是光滑二维链的边界。再用第一卷定理14.5.5的链上 Stokes 公式，对实部与虚部分别计算，得到
\[
	\int_{\partial c}\alpha=\int_c\mathrm{d}\alpha=0.
\]
于是可定义 \(\alpha\) 的 \(a\)-周期、\(b\)-周期
\[
	A_i(\alpha)=\int_{a_i}\alpha,
	\qquad B_i(\alpha)=\int_{b_i}\alpha.
\]
它们合起来确定同态 \(H_1(X;\mathbb Z)\to\mathbb C\)，称为 \(\alpha\) 的\term[zhouqi]{周期}[period]同态。

\ProseParagraphBreak

全纯微分 \(\omega=f(z)\,\mathrm{d}z\) 是闭形式，因为
\[
	\mathrm{d}(f\,\mathrm{d}z)
	=\frac{\partial f}{\partial\bar z}\,\mathrm{d}\bar z\wedge\mathrm{d}z=0.
\]
由\cref{g4:thm:0304:genus-comparison}，\(H^0(X,K_X)\) 的复维数为 \(g\)。暂取任意一组全纯微分基 \(\eta_1,\ldots,\eta_g\)，写
\[
	A_{ij}=\int_{a_i}\eta_j,
	\qquad B_{ij}=\int_{b_i}\eta_j.
\]
矩阵 \(A,B\) 的行对应闭路，列对应微分；这一约定决定后面归一化时矩阵相乘的次序。

\subsection{Riemann 双线性关系}
\label{g4:subsec:0401:02}

两组周期之间的关系来自一个二维积分。把曲面切成盘后，闭形式有单值原函数，可以对“原函数乘另一闭形式”使用 Stokes 公式。重新粘合时，原函数在配对边上的跳跃正好是周期。

\begin{theorem}
	\label{g4:thm:0401:bilinear-periods}
	设 \(X\) 是紧连通黎曼曲面，取\cref{g4:prop:0401:symplectic-basis}中的同调辛基。对于任意光滑复值闭一形式 \(\alpha,\beta\)，有
	\begin{equation}\label{g4:eq:0401:bilinear-periods}
		\int_X\alpha\wedge\beta
		=\sum_{i=1}^g
		\bigl(A_i(\alpha)B_i(\beta)-B_i(\alpha)A_i(\beta)\bigr).
	\end{equation}
\end{theorem}
\begin{proof}
	先设 \(g\ge1\)，沿式\eqref{g4:eq:0401:polygon-word}切开。闭形式 \(\alpha\) 在盘 \(P\) 上有原函数 \(F\)。具体地，先在内部固定一点并按道路积分定义 \(F\)；任意两条道路之差在开盘中零同调，前述链上 Stokes 公式说明积分相同。局部原函数又说明 \(F\) 光滑，并可按各边和角的局部片延拓到边界。最后减去常数，使一个选定顶点副本处的值为零。因此
	\[
		\mathrm{d}(F\beta)=\mathrm{d}F\wedge\beta+F\,\mathrm{d}\beta
		=\alpha\wedge\beta.
	\]
	切口是一维有限曲线网，面积为零。第一卷定理14.4.5的带角 Stokes 公式遂给出
	\[
		\int_X\alpha\wedge\beta=\int_P\mathrm{d}(F\beta)
		=\int_{\partial P}F\beta.
	\]
	也可将 \(P\) 分成有限个带角块应用该公式，内部公共边的积分两两抵消。

	下面逐对计算边界项，同时固定其符号。记第 \(i\) 个四边块起点的原函数值为 \(C_i\)，并沿 \(a_i,b_i\) 的正向参数记
	\[
		f_a(t)=\int_{a_i|_{[0,t]}}\alpha,
		\qquad f_b(t)=\int_{b_i|_{[0,t]}}\alpha.
	\]
	用 \(a_i^+,b_i^+\) 表示在边界中正向出现的边，用 \(a_i^-,b_i^-\) 表示其配对边；表中负边的点仍按对应正边的参数 \(t\) 标记。沿正向边界累加积分，得到
	\begin{equation}\label{g4:eq:0401:edge-jumps}
		\begin{array}{c|cc}
		 & \text{正边上的值} & \text{负边匹配点的值}\\ \hline
		a_i & C_i+f_a(t) & C_i+B_i(\alpha)+f_a(t)\\
		b_i & C_i+A_i(\alpha)+f_b(t) & C_i+f_b(t).
		\end{array}
	\end{equation}
	例如沿 \(a_i b_i\) 走完时值为 \(C_i+A_i(\alpha)+B_i(\alpha)\)；再沿 \(a_i^{-1}\) 走到与参数 \(t\) 对应的点，须减去 \(A_i(\alpha)-f_a(t)\)，即得表中第一行右项。第二行同样得到。每个四边块的总增量为零，故下一个 \(C_{i+1}\) 与 \(C_i\) 相同。

	配对边上 \(\beta\) 是同一个形式，而边界方向相反。因此 \(a_i\) 这对边贡献
	\[
		\int_{a_i}\bigl(F|_{a_i^+}-F|_{a_i^-}\bigr)\beta
		=-B_i(\alpha)A_i(\beta),
	\]
	\(b_i\) 这对边贡献 \(A_i(\alpha)B_i(\beta)\)。求和就是式\eqref{g4:eq:0401:bilinear-periods}。

	当 \(g=0\) 时，\(H_1(X;\mathbb Z)=0\)。闭形式的道路积分给出全局原函数 \(F\)，于是 \(\alpha\wedge\beta=\mathrm{d}(F\beta)\)。无边界 Stokes 公式说明积分为零，与右边的空和一致。
\end{proof}

公式也适用于任意同调辛基。确实，按先 \(a\)、后 \(b\) 排列，令
\(J=\left(\begin{smallmatrix}0&I_g\\-I_g&0\end{smallmatrix}\right)\)。
若新基由整数矩阵 \(M\) 给出，保持相交数就是 \(MJM^{\mathsf T}=J\)，它等价于 \(M^{\mathsf T}JM=J\)。周期列向量同时乘以 \(M\)，而式\eqref{g4:eq:0401:bilinear-periods}的右侧是 \(v_\alpha^{\mathsf T}Jv_\beta\)，故其数值不变。

\begin{theorem}[Riemann 双线性关系]
	\label{g4:thm:0401:riemann-relations}
	设 \(X\) 是紧连通黎曼曲面，\(a_i,b_i\) 是上述同调辛基。任意全纯微分 \(\omega,\eta\) 满足
	\begin{equation}\label{g4:eq:0401:first-riemann}
		\sum_{i=1}^g
		\bigl(A_i(\omega)B_i(\eta)-B_i(\omega)A_i(\eta)\bigr)=0.
	\end{equation}
	若 \(\omega\ne0\)，则
	\begin{equation}\label{g4:eq:0401:second-riemann}
		\frac{\mathrm{i}}2\int_X\omega\wedge\overline\omega
		=\operatorname{Im}\sum_{i=1}^g
		\overline{A_i(\omega)}B_i(\omega)>0.
	\end{equation}
\end{theorem}
\begin{proof}
	在局部坐标中，\(\omega\wedge\eta=f(z)h(z)\,\mathrm{d}z\wedge\mathrm{d}z=0\)，故第一式由\cref{g4:thm:0401:bilinear-periods}得到。第二式使用 \(\beta=\overline\omega\)。因为沿实参数道路积分与复共轭交换，右边为
	\[
		\frac{\mathrm{i}}2\sum_i
		\left(A_i(\omega)\overline{B_i(\omega)}
		-B_i(\omega)\overline{A_i(\omega)}\right)
		=\operatorname{Im}\sum_i\overline{A_i(\omega)}B_i(\omega).
	\]
	另一方面，写 \(z=x+\mathrm{i}y\) 及 \(\omega=f(z)\,\mathrm{d}z\)，则
	\[
		\frac{\mathrm{i}}2\omega\wedge\overline\omega
		=|f(z)|^2\,\mathrm{d}x\wedge\mathrm{d}y.
	\]
	它处处非负，并在 \(\omega\) 非零的一点附近严格为正。紧性保证积分有限，非零性保证积分严格为正。
\end{proof}

第一关系来自两个 \((1,0)\) 形式的楔积为零，第二关系来自 \((1,0)\) 形式与其共轭的楔积给出正面积。这两种局部计算经过切割积分，分别变成周期的代数约束。

\subsection{归一化全纯微分}
\label{g4:subsec:0401:03}

第二双线性关系首先排除一种可能：非零全纯微分不能让全部 \(a\)-周期同时为零。因此，\(a\)-周期可以用作全纯微分空间的坐标。

\begin{theorem}
	\label{g4:thm:0401:normalized-periods}
	设 \(X\) 是亏格 \(g\ge1\) 的紧连通黎曼曲面，并固定同调辛基 \(a_i,b_i\)。存在唯一的全纯微分基 \(\omega_1,\ldots,\omega_g\)，满足
	\begin{equation}\label{g4:eq:0401:normalized-a}
		\int_{a_i}\omega_j=\delta_{ij}.
	\end{equation}
	称它为相对于所取 \(a\)-周期的\term[guiyihuaquanchunweifen]{归一化全纯微分}[normalized holomorphic differentials]。其周期矩阵
	\begin{equation}\label{g4:eq:0401:tau-definition}
		\tau_{ij}\coloneqq\int_{b_i}\omega_j
	\end{equation}
	满足
	\begin{equation}\label{g4:eq:0401:tau-properties}
		\tau^{\mathsf T}=\tau,
		\qquad Y\coloneqq\operatorname{Im}\tau>0.
	\end{equation}
	这里 \(Y>0\) 指实对称矩阵 \(Y\) 正定。
\end{theorem}
\begin{proof}
	考虑线性映射
	\[
		H^0(X,K_X)\longrightarrow\mathbb C^g,
		\qquad\omega\longmapsto(A_1(\omega),\ldots,A_g(\omega))^{\mathsf T}.
	\]
	若其像为零，式\eqref{g4:eq:0401:second-riemann}迫使 \(\omega=0\)，所以它单射。由\cref{g4:thm:0304:genus-comparison}，源与靶的复维数同为 \(g\)，故它是同构。各标准基向量的唯一原像就是 \(\omega_j\)。

	在第一关系中取 \(\omega_j,\omega_k\)，得到 \(\tau_{jk}-\tau_{kj}=0\)。若 \(c=(c_1,\ldots,c_g)^{\mathsf T}\ne0\)，令 \(\omega=\sum_jc_j\omega_j\)，则其周期列向量为
	\[
		A(\omega)=c,\qquad B(\omega)=\tau c.
	\]
	所以第二关系给出
	\[
		0<\frac{\mathrm{i}}2\int_X\omega\wedge\overline\omega
		=\operatorname{Im}(\overline c^{\mathsf T}\tau c)
		=\overline c^{\mathsf T}Yc.
	\]
	最后一个等式使用 \(\operatorname{Re}\tau\) 与 \(Y\) 都是实对称矩阵，故 \(\overline c^{\mathsf T}(\operatorname{Re}\tau)c\) 及 \(\overline c^{\mathsf T}Yc\) 均为实数。尤其对非零实向量成立，证明 \(Y\) 正定。
\end{proof}

对于前面任取的微分基 \(\eta_j\)，矩阵 \(A\) 因而可逆，归一化明确写成
\[
	\omega_j=\sum_k\eta_k(A^{-1})_{kj},
	\qquad \tau=BA^{-1}.
\]
归一化改变了微分的坐标，没有改变全纯微分空间。\(\tau\) 仍依赖所选同调辛基；下一节将用一个不依赖坐标的商空间解释这种依赖。

\begin{example}
	\label{g4:exm:0401:torus-periods}
	设 \(X=\mathbb C/(\mathbb Z\lambda_1+\mathbb Z\lambda_2)\)，其中 \(\lambda_1,\lambda_2\) 实线性无关，并按 \(\operatorname{Im}(\lambda_2/\lambda_1)>0\) 排列。沿 \(\lambda_1,\lambda_2\) 的线段在商中给出 \(a,b\)，且 \(a\cdot b=1\)。由\cref{g4:prop:0201:torus}的构造，\(\mathrm{d}z\) 在格平移下不变，故下降到 \(X\)。归一化微分与周期矩阵为
	\[
		\omega=\frac{\mathrm{d}z}{\lambda_1},
		\qquad \tau=\frac{\lambda_2}{\lambda_1}.
	\]
	第二关系在此就是正向基本平行四边形的面积为正。交换 \(a,b\) 而不改变其中一条的方向，会使相交数变成 \(-1\)，因而不能保留原来的辛基约定。
\end{example}

\subsection{周期格的离散性与正性}
\label{g4:subsec:0401:04}

把同一闭路上全部归一化微分的积分排成一个列向量，得到
\[
	\operatorname{Per}:H_1(X;\mathbb Z)\longrightarrow\mathbb C^g,
	\qquad \gamma\longmapsto
	\left(\int_\gamma\omega_1,\ldots,\int_\gamma\omega_g\right)^{\mathsf T}.
\]
\(a_i\) 的像是第 \(i\) 个标准基向量；\(b_i\) 的像是 \(\tau\) 的第 \(i\) 行转置。因为 \(\tau\) 对称，这个转置正好是 \(\tau\) 的第 \(i\) 列。于是周期群为
\begin{equation}\label{g4:eq:0401:period-lattice}
	\Lambda\coloneqq\operatorname{Per}\bigl(H_1(X;\mathbb Z)\bigr)
	=\mathbb Z^g+\tau\mathbb Z^g.
\end{equation}

\begin{proposition}
	\label{g4:prop:0401:period-lattice}
	设 \(X\) 是亏格 \(g\ge1\) 的紧连通黎曼曲面，\(\tau\) 是其归一化周期矩阵。群 \(\Lambda=\mathbb Z^g+\tau\mathbb Z^g\) 是 \(\mathbb C^g\) 中秩 \(2g\) 的离散子群，其实线性张成空间为 \(\mathbb C^g\)，且 \(\mathbb C^g/\Lambda\) 紧。特别地，\(\operatorname{Per}\) 为单射。
\end{proposition}
\begin{proof}
	写 \(\tau=U+\mathrm{i}Y\)。正定性保证 \(Y\) 可逆，所以实线性映射
	\[
		\Phi:\mathbb R^g\oplus\mathbb R^g\longrightarrow\mathbb C^g,
		\qquad(s,t)\longmapsto s+\tau t
	\]
	是同构：给定 \(z\)，其逆明确为
	\[
		t=Y^{-1}\operatorname{Im}z,
		\qquad s=\operatorname{Re}z-UY^{-1}\operatorname{Im}z.
	\]
	因此存在 \(c>0\)，使所有实向量满足
	\[
		\|s+\tau t\|\ge c\sqrt{\|s\|^2+\|t\|^2}.
	\]
	非零整数向量的右端至少为 \(c\)，故不同格点的距离至少为 \(c\)。这证明离散性，同时说明 \(2g\) 个生成向量实线性无关，且整数表示唯一。

	任意 \((s,t)\) 减去一个整数向量后落在 \([0,1]^{2g}\) 中，所以紧平行多面体 \(\Phi([0,1]^{2g})\) 的像覆盖整个商空间，商因而紧。最后，\(a_i,b_i\) 是整数同调基，且它们的周期像整数线性无关，故周期同态单射。
\end{proof}

实向量空间中由一组实基的整数线性组合构成的子群称为满秩\term[ge]{格}[lattice]。这里正性既保证其离散，又提供一个以后有用的正 Hermitian 形式。沿用第三章对第一变量线性的约定，令
\begin{equation}\label{g4:eq:0401:polarization-form}
	H(z,w)=z^{\mathsf T}Y^{-1}\overline w,
	\qquad E(z,w)=-\operatorname{Im}H(z,w).
\end{equation}
\(H\) 对第一变量线性、对第二变量共轭线性且正定，并有 \(E(v,\mathrm{i}v)=H(v,v)>0\)（\(v\ne0\)）。若 \(z=m+\tau n\)、\(w=m'+\tau n'\)，其中 \(m,n,m',n'\in\mathbb Z^g\)，直接展开并用 \(U^{\mathsf T}=U\) 得
\begin{equation}\label{g4:eq:0401:integral-alternating-form}
	E(z,w)=m^{\mathsf T}n'-n^{\mathsf T}m'.
\end{equation}
因此 \(E\) 在格上取整数值，其矩阵为 \(\left(\begin{smallmatrix}0&I_g\\-I_g&0\end{smallmatrix}\right)\)，行列式为 \(1\)。它正好保留了原曲面的相交配对；第4节讨论主极化时将回到这个形式。

\ProseParagraphBreak

当 \(g=0\) 时，全纯微分空间为零，约定 \(\mathbb C^0=\{0\}\)、\(\Lambda=\{0\}\)，相应商空间只有一点。对于正亏格曲面，周期则把全部整数一维同调嵌入一个复 \(g\) 维空间。这是下一节构造 Jacobi 环面的基础。
% 纲要条目排为编号小节标题，并进入目录与书签；说明段仍保留于纲要注释中。
% !TeX root = ../../Book4.tex
% 纲要标题：### 4.2 Jacobi 簇和 Abel–Jacobi 映射
% 纲要位置：几何学计划纲要.md，第 3599 行。
% 正文状态：初稿；尚未正式审核。

% 纲要内容开始
% 
% * 复向量空间模周期格
% * 路径积分与基点依赖
% * 对称积上的映射
% * 除子类的几何解释
% 
% 纲要内容结束

\section{Jacobi 簇和 Abel–Jacobi 映射}
\label{g4:sec:0402}

固定起点后，把全纯微分从起点积到 \(p\)，会得到随 \(p\) 变化的一组复数。改变道路使这些复数增加周期，因此它们在上一节的周期格取商后成为单值。这既构造出一个复环面，也把曲面本身以及曲面上的点组映入这个环面。为了允许点组中的点相互重合，我们还要为无序点组建立真正的复坐标。

\subsection{复向量空间模周期格}
\label{g4:subsec:0402:01}

本节出现复维数大于一的空间。这里的复 \(n\) 维流形是 Hausdorff、第二可数空间，带有到 \(\mathbb C^n\) 开集的坐标图，换坐标及其逆都全纯。可把坐标函数的全纯性理解为局部由收敛的多变量复幂级数表示；下文所用的平移、线性映射、多项式及含参数的积分都会直接满足这一条件。复维数一时，这与黎曼曲面的定义一致。

\begin{proposition}
	\label{g4:prop:0402:complex-torus}
	设 \(V\) 是复 \(g\) 维向量空间，\(\Lambda\subset V\) 是实秩 \(2g\) 的满秩格。商 \(V/\Lambda\) 有唯一使投影局部双全纯的复流形结构。它紧、连通，复维数为 \(g\)，加法与取负都是全纯映射。
\end{proposition}
\begin{proof}
	选择一个 Euclid 范数。格的离散性给出 \(\varepsilon>0\)，使不同格点的距离至少为 \(3\varepsilon\)。因此半径小于 \(\varepsilon\) 的球与它的任何非零格平移不交，投影 \(\pi:V\to V/\Lambda\) 在该球上单射。商映射是开映射，因为
	\[
		\pi^{-1}(\pi(U))=\bigcup_{\lambda\in\Lambda}(U+\lambda)
	\]
	对开集 \(U\) 仍为开集，所以这些局部投影都是到开集的同胚。

	先核对商的分离性。若 \(z,w\) 不属于同一格轨道，则 \(z-w\notin\Lambda\)。一个有界集只含有限个格点，故 \(z-w\) 到 \(\Lambda\) 的距离严格为正。取充分小的球 \(B_z,B_w\)，使任何 \(B_z\) 与 \(B_w+\lambda\) 都不交；它们在商中的像就是两个不交邻域。因此商 Hausdorff。\(V\) 的可数开基在开映射 \(\pi\) 下的像仍构成可数开基，所以商第二可数。

	用前述小球给商定义复坐标。两张图的换坐标在每个交叠分支上相差一个固定格向量，故为全纯平移。这同时给出存在性与唯一性。格的一组实基所张成的闭基本平行多面体紧，其像覆盖商，故商紧；\(V\) 连通，所以商也连通。加法在局部提升后为 \((z,w)\mapsto z+w\)，取负为 \(z\mapsto-z\)，均为全纯映射。
\end{proof}

这样的商称为\term[fuhuanmian]{复环面}[complex torus]。下面用全纯微分产生的特定周期格来构造它。

\begin{definition}
	\label{g4:def:0402:jacobian}
	设 \(X\) 是亏格 \(g\) 的紧连通黎曼曲面。固定同调辛基和归一化全纯微分，令 \(\tau\) 为周期矩阵、\(\Lambda=\mathbb Z^g+\tau\mathbb Z^g\)。其\term[jacobicu]{Jacobi 簇}[Jacobian variety]的复解析表示为
	\[
		J(X)\coloneqq\mathbb C^g/\Lambda.
	\]
	亏格零时约定 \(J(X)\) 为一点。
\end{definition}

\cref{g4:prop:0401:period-lattice,g4:prop:0402:complex-torus}保证这个构造确为紧复环面，复维数为 \(g\)。其上的额外正性由式\eqref{g4:eq:0401:polarization-form}记录，第4节再讨论它与主极化的关系。

\ProseParagraphBreak

这个对象也有不选择基的定义。记 \(V=H^0(X,K_X)^*\) 为全纯微分空间的复线性对偶。每个整数同调类 \(\gamma\) 定义一个泛函
\[
	\ell_\gamma:H^0(X,K_X)\longrightarrow\mathbb C,
	\qquad \ell_\gamma(\omega)=\int_\gamma\omega.
\]
取商
\begin{equation}\label{g4:eq:0402:intrinsic-jacobian}
	J(X)=H^0(X,K_X)^*/\{\ell_\gamma:\gamma\in H_1(X;\mathbb Z)\}.
\end{equation}
在归一化基下，泛函 \(\ell\) 的坐标是列向量
\((\ell(\omega_1),\ldots,\ell(\omega_g))^{\mathsf T}\)，周期子群恰变成 \(\Lambda\)。因此式\eqref{g4:eq:0402:intrinsic-jacobian}与前一定义一致。改变微分基只是对 \(V\) 作复线性换坐标；改变同调基只改变同一周期子群的整数生成元。两种改变均给出同一个内在复环面，而不是另造一个空间。

\subsection{路径积分与基点依赖}
\label{g4:subsec:0402:02}

设 \(q\in X\) 为基点。曲面连通且局部道路连通，任意 \(p\in X\) 都可从 \(q\) 通过有限条坐标圆盘中的光滑弧到达。沿这样一条分片光滑道路 \(\gamma_{qp}\)，定义积分列向量
\[
	\left(\int_{\gamma_{qp}}\omega_1,\ldots,
	\int_{\gamma_{qp}}\omega_g\right)^{\mathsf T}.
\]
这个向量通常依赖道路；它的商类具有下面的性质。

\begin{proposition}
	\label{g4:prop:0402:abel-jacobi}
	设 \(X\) 是紧连通黎曼曲面，\(q\in X\)。公式
	\begin{equation}\label{g4:eq:0402:abel-jacobi}
		A_q(p)\coloneqq
		\left[\left(\int_q^p\omega_1,\ldots,\int_q^p\omega_g\right)^{\mathsf T}\right]
		\in J(X)
	\end{equation}
	与所取道路无关，并定义全纯映射，称为基于 \(q\) 的\term[abeljacobiyingshe]{Abel--Jacobi 映射}[Abel--Jacobi map]。它满足 \(A_q(q)=0\)。若改用基点 \(q'\)，则
	\begin{equation}\label{g4:eq:0402:basepoint-change}
		A_{q'}(p)=A_q(p)-A_q(q').
	\end{equation}
	以平移识别 \(T_{A_q(p)}J(X)\) 与 \(\mathbb C^g\) 后，其微分为
	\begin{equation}\label{g4:eq:0402:abel-tangent}
		\mathrm{d}(A_q)_p(v)
		=\bigl(\omega_1(v),\ldots,\omega_g(v)\bigr)^{\mathsf T}.
	\end{equation}
\end{proposition}
\begin{proof}
	两条从 \(q\) 到 \(p\) 的道路之差是一条整数一圈；各全纯微分闭，积分之差只依赖这个圈的同调类。因此两个积分向量相差 \(\Lambda\) 中的元素，在商中相等。

	在 \(p_0\) 附近取单连通坐标圆盘 \(U\)，先固定从 \(q\) 到 \(p_0\) 的道路，再在 \(U\) 内积分。每个 \(\omega_j\) 在 \(U\) 上有全纯原函数 \(F_j\)，于是 \(A_q|_U\) 有一个到 \(\mathbb C^g\) 的提升
	\[
		p\longmapsto c+
		\bigl(F_1(p)-F_1(p_0),\ldots,F_g(p)-F_g(p_0)\bigr)^{\mathsf T}.
	\]
	它全纯，求导便得到式\eqref{g4:eq:0402:abel-tangent}。不同提升相差常数格向量，所以微分不受影响。

	若先从 \(q\) 走到 \(q'\)，再从 \(q'\) 走到 \(p\)，积分的可加性给出
	\(A_q(p)=A_q(q')+A_{q'}(p)\)，即得基点公式。常值道路给出 \(A_q(q)=0\)。
\end{proof}

也可以直接按式\eqref{g4:eq:0402:intrinsic-jacobian}，把 \(A_q(p)\) 看成泛函 \(\omega\mapsto\int_q^p\omega\) 的商类。这样，映射的定义只需基点，不需选择微分基。基点变化只在目标中作平移，故不会改变映射的局部微分性质。

\begin{example}
	\label{g4:exm:0402:torus-abel}
	沿用\cref{g4:exm:0401:torus-periods}的复环面，基点取 \(q=[0]\)。归一化微分为 \(\mathrm{d}z/\lambda_1\)，故
	\[
		A_q:\mathbb C/(\mathbb Z\lambda_1+\mathbb Z\lambda_2)
		\longrightarrow\mathbb C/(\mathbb Z+\tau\mathbb Z),
		\qquad [z]\longmapsto[z/\lambda_1].
	\]
	它是复 Lie 群同构，逆为 \([w]\mapsto[\lambda_1w]\)。在这个模型中，沿不同道路绕行造成的积分差恰好被商中的格吸收。
\end{example}

\subsection{对称积上的映射}
\label{g4:subsec:0402:03}

若有 \(d\) 个点，应把它们的积分相加。这个和不记录点的次序，所以其定义域应是无序点组，而不是有序积 \(X^d\)。在两点重合时，交换两个点具有不动点；一般的有限群商可能在这里出现奇点。曲线的特殊之处是，无序点组可改用一元多项式的系数记录，从而仍然得到光滑的复流形。

\begin{definition}
	\label{g4:def:0402:symmetric-product}
	设 \(X\) 是黎曼曲面，\(d\ge1\)。对称群 \(\mathfrak S_d\) 通过置换因子作用于 \(X^d\)，商
	\[
		X^{(d)}\coloneqq X^d/\mathfrak S_d
	\]
	赋予商拓扑，称为 \(X\) 的 \(d\) 次\term[duichengji]{对称积}[symmetric product]。它的点可写为有效除子
	\(D=p_1+\cdots+p_d=\sum_{\nu=1}^r m_\nu p_\nu\)，其中各 \(p_\nu\) 不同，\(m_\nu\ge1\)，\(\sum_\nu m_\nu=d\)。并约定 \(X^{(0)}\) 为一点。
\end{definition}

\begin{theorem}
	\label{g4:thm:0402:symmetric-product}
	设 \(X\) 是黎曼曲面，\(d\ge1\)。空间 \(X^{(d)}\) 有自然的复 \(d\) 维流形结构；在 \(D=\sum_\nu m_\nu p_\nu\) 附近，取各 \(p_\nu\) 的互不相交坐标圆盘，以每组 \(m_\nu\) 个坐标值的初等对称函数为复坐标。商映射 \(X^d\to X^{(d)}\) 全纯。若 \(X\) 紧，则 \(X^{(d)}\) 紧。
\end{theorem}
\begin{proof}
	先看一个圆盘 \(D_r=\{z:|z|<r\}\) 中的 \(m\) 个点。将无序点组 \(\{z_1,\ldots,z_m\}\) 送到系数 \((\sigma_1,\ldots,\sigma_m)\)，其中
	\begin{equation}\label{g4:eq:0402:symmetric-polynomial}
		P_\sigma(t)=\prod_{\ell=1}^m(t-z_\ell)
		=t^m-\sigma_1t^{m-1}+\cdots+(-1)^m\sigma_m.
	\end{equation}
	它连续，并由复多项式的根及其重数唯一决定，因此是到集合
	\[
		\Omega_{m,r}=\{\sigma\in\mathbb C^m:
		P_\sigma\text{ 的全部根都在 }D_r\}
	\]
	的双射。

	集合 \(\Omega_{m,r}\) 是开集。给定其中一点，选半径 \(R<r\)，使所有根都在 \(|t|<R\)。原多项式在 \(|t|=R\) 上有严格正的最小模；系数充分接近时，多项式之差在该圆周上的模小于这一最小值。\cref{g4:thm:0102:rouche}的 Rouché 定理说明新多项式仍有全部 \(m\) 个根位于 \(|t|<R\)。逆映射也连续：在原多项式的各不同根周围取任意小的互不相交圆盘；再用 Rouché 定理，充分小的系数扰动在每张盘内保留原来的根数，按重数计。这正表示根组成的无序点组趋近原点组。因此 \(D_r^{(m)}\to\Omega_{m,r}\) 为同胚。

	现在在 \(D=\sum_\nu m_\nu p_\nu\) 的各支撑点周围取互不相交坐标圆盘 \(U_\nu\)，坐标 \(z_\nu(p_\nu)=0\)。恰在每个 \(U_\nu\) 中含 \(m_\nu\) 个点的无序点组构成 \(D\) 的开邻域，且同胚于 \(\prod_\nu U_\nu^{(m_\nu)}\)。前面的系数坐标把它映到 \(\mathbb C^{m_1}\times\cdots\times\mathbb C^{m_r}\) 的开集。

	还要证明换坐标全纯，尤其在重根处不能通过逐个选根来证明。设 \(w=\varphi(z)\) 是两张曲面坐标之间的全纯变换。围绕当前点组的各不同根，取位于两张坐标交集内的小圆周，合起来记为 \(\Gamma\)，全部正向。系数在一个小邻域内变化时，没有根穿过 \(\Gamma\)。对于被这些圆周包围的根，留数定理给出幂和公式
	\begin{equation}\label{g4:eq:0402:trace-integral}
		\sum_{\ell}\varphi(z_\ell)^k
		=\frac{1}{2\pi\mathrm{i}}\int_\Gamma
		\varphi(t)^k\,\frac{P_\sigma'(t)}{P_\sigma(t)}\,\mathrm{d}t,
		\qquad k\ge1.
	\end{equation}
	重根 \(z_\ell\) 的留数按其重数计入，所以该式也覆盖重合情形。右侧对 \(\sigma\) 全纯：分母在固定紧圆周上离零有正距离，对附近系数可把其倒数展开为一致收敛的几何级数；逐项积分给出系数变量的收敛幂级数。

	记新坐标根值的幂和为 \(S_k\)，其初等对称函数为 \(e_k\)，则
	\[
		e_0=1,\qquad
		k e_k=\sum_{\ell=1}^k(-1)^{\ell-1}e_{k-\ell}S_\ell.
	\]
	这个恒等式可比较 \(\prod_\ell(1+w_\ell u)\) 的对数导数两边的 \(u\) 次幂得到。因此各 \(e_k\) 也是旧系数的全纯函数。若两张点组图把根分到不同的圆盘中，就在当前各支撑点处取共同的更小圆盘，用式\eqref{g4:eq:0402:trace-integral}分别选出每组根，再相加得到所需幂和。同一论证反向使用，说明逆换坐标也全纯。

	商空间的 Hausdorff 性来自有限群作用：两个不同有限轨道可在 Hausdorff 空间 \(X^d\) 中用不交开集分离，再取有限个群作用后的交使邻域不变。商映射是开映射，故 \(X^d\) 的可数开基给出商的可数基。以上图册因此构成复流形结构。商映射在这些坐标中就是初等对称多项式，故全纯；若 \(X\) 紧，商作为 \(X^d\) 的连续像也紧。
\end{proof}

\begin{example}
	\label{g4:exm:0402:double-point}
	在 \(2p\) 附近取坐标 \(z(p)=0\)，则 \(X^{(2)}\) 的坐标为
	\[
		\sigma_1=z_1+z_2,\qquad \sigma_2=z_1z_2.
	\]
	重合点组对应判别式曲线 \(\sigma_1^2-4\sigma_2=0\)。商映射的 Jacobian 行列式为 \(z_1-z_2\)，在对角线上为零，但目标的 \((\sigma_1,\sigma_2)\) 仍是普通复坐标。这里退化的是从有序点组到无序点组的映射，而非无序点组空间本身。
\end{example}

\begin{proposition}
	\label{g4:prop:0402:symmetric-abel}
	设 \(X\) 是紧连通黎曼曲面，\(q\in X\)，\(d\ge0\)。公式
	\begin{equation}\label{g4:eq:0402:symmetric-abel}
		A_q^{(d)}:X^{(d)}\longrightarrow J(X),
		\qquad p_1+\cdots+p_d\longmapsto\sum_{\ell=1}^d A_q(p_\ell)
	\end{equation}
	定义全纯映射，称为对称积上的 Abel 映射。其基点变化公式为
	\[
		A_{q'}^{(d)}(D)=A_q^{(d)}(D)-dA_q(q').
	\]
	当 \(d=0\) 时取值为 \(0\)。
\end{proposition}
\begin{proof}
	求和不依赖点的排列，故得到连续的商映射。基点公式由式\eqref{g4:eq:0402:basepoint-change}逐项相加得到。全纯性要在刚建立的系数坐标中检查。

	固定 \(D=\sum_\nu m_\nu p_\nu\)，在各坐标圆盘 \(U_\nu\) 上取微分 \(\omega_j\) 的全纯原函数 \(F_{j,\nu}\)，并固定从 \(q\) 到 \(p_\nu\) 的道路。附近点组的积分和有一个局部提升，其第 \(j\) 个分量为常数加
	\[
		\sum_{\nu=1}^r\sum_{\ell=1}^{m_\nu}
		F_{j,\nu}(z_{\nu,\ell}).
	\]
	以 \(F_{j,\nu}\) 代替式\eqref{g4:eq:0402:trace-integral}中的 \(\varphi^k\)，得到
	\[
		\sum_{\ell=1}^{m_\nu}F_{j,\nu}(z_{\nu,\ell})
		=\frac{1}{2\pi\mathrm{i}}\int_{\Gamma_\nu}
		F_{j,\nu}(t)\frac{P_{\nu,\sigma}'(t)}
		{P_{\nu,\sigma}(t)}\,\mathrm{d}t.
	\]
	这里 \(\Gamma_\nu\) 是固定在 \(U_\nu\) 内、围住这组根的小圆周。右侧按前一证明对系数全纯，所以局部提升全纯。这也证明了映射在重复点处的全纯性。
\end{proof}

在各点不同的点组上，微分就是各点求值的和。重复点则把求值扩展为有限阶喷射。下面的公式将使第四节能够用微分的消失阶数判断 Abel 映射的秩。

\begin{proposition}
	\label{g4:prop:0402:abel-differential}
	设 \(X\) 是亏格 \(g\) 的紧连通黎曼曲面，\(q\in X\)，\(d\ge1\)，且 \(D=\sum_{\nu=1}^r m_\nu p_\nu\in X^{(d)}\)。通过
	\[
		T_{A_q^{(d)}(D)}^*J(X)\cong H^0(X,K_X)
	\]
	的自然识别，微分 \(\mathrm{d}(A_q^{(d)})_D\) 的转置之核恰为 \(H^0(X,K_X(-D))\)。因此
	\begin{equation}\label{g4:eq:0402:abel-rank}
		\operatorname{rank}_{\mathbb C}\mathrm{d}(A_q^{(d)})_D
		=g-h^0(K_X(-D)).
	\end{equation}
	具体地，在 \(p_\nu\) 附近取坐标 \(z_\nu(p_\nu)=0\)，若
	\(\omega_j=\bigl(\sum_{k\ge0}a_{j,\nu,k}z_\nu^k\bigr)\,\mathrm{d}z_\nu\)，则以初等对称函数 \(\sigma_{\nu,k}\) 为坐标，有
	\begin{equation}\label{g4:eq:0402:abel-jet-matrix}
		\frac{\partial(A_q^{(d)})_j}{\partial\sigma_{\nu,k}}(D)
		=(-1)^{k-1}a_{j,\nu,k-1},
		\qquad 1\le k\le m_\nu.
	\end{equation}
	左侧在任意局部提升中计算，结果相同。
\end{proposition}
\begin{proof}
	只需在一组重合点 \(mp\) 附近计算，其他组的贡献相加。记局部原函数为 \(F_j\)，使 \(F_j'(z)=\sum_{k\ge0}a_{j,k}z^k\)。在重合点，式\eqref{g4:eq:0402:symmetric-polynomial}中的多项式为 \(P_0(t)=t^m\)。沿任意系数方向的变化写成
	\[
		\delta P(t)=\sum_{k=1}^m(-1)^k\delta\sigma_k\,t^{m-k}.
	\]
	对上一证明的积分求导，并沿闭圆周分部积分，得到
	\[
		\begin{aligned}
		\delta\left(\sum_{\ell=1}^mF_j(z_\ell)\right)
		&=\frac{1}{2\pi\mathrm{i}}\int_\Gamma
		F_j(t)\,\mathrm{d}\left(\frac{\delta P(t)}{P_0(t)}\right)\\
		&=-\frac{1}{2\pi\mathrm{i}}\int_\Gamma
		F_j'(t)\frac{\delta P(t)}{t^m}\,\mathrm{d}t\\
		&=\sum_{k=1}^m(-1)^{k-1}a_{j,k-1}\,\delta\sigma_k.
		\end{aligned}
	\]
	求导与积分交换由固定圆周上的一致幂级数保证；分部积分中的全微分沿闭路积分为零。由留数取出 \(t^{-1}\) 系数，得到最后一行，亦即式\eqref{g4:eq:0402:abel-jet-matrix}。

	在内在表示 \(J(X)=H^0(K_X)^*/\operatorname{Per}(H_1)\) 下，平移把其切空间识别为 \(H^0(K_X)^*\)，所以余切空间自然是 \(H^0(K_X)\)。一个余切向量 \(\omega=\sum_jc_j\omega_j\) 消去微分的像，当且仅当在每个 \(p_\nu\) 处的前 \(m_\nu\) 个 Taylor 系数均为零。这等价于 \(\operatorname{ord}_{p_\nu}\omega\ge m_\nu\)，即 \(\omega\in H^0(K_X(-D))\)。最后用有限维线性代数中转置的秩等于原映射的秩，得到式\eqref{g4:eq:0402:abel-rank}。
\end{proof}

\subsection{除子类的几何解释}
\label{g4:subsec:0402:04}

有效除子记录点组，两个同次数点组之差则属于
\[
	\operatorname{Div}^0(X)=\left\{\sum_p n_p[p]:
	n_p\in\mathbb Z,\ \text{仅有限项非零},\ \sum_p n_p=0\right\}.
\]
对这样的差，基点变化产生的常数恰好全部抵消。

\begin{proposition}
	\label{g4:prop:0402:divisor-integral}
	设 \(X\) 是紧连通黎曼曲面。公式
	\begin{equation}\label{g4:eq:0402:divisor-integral}
		\operatorname{AJ}:\operatorname{Div}^0(X)\longrightarrow J(X),
		\qquad
		\operatorname{AJ}\left(\sum_p n_p[p]\right)
		=\sum_p n_pA_q(p)
	\end{equation}
	定义与基点 \(q\) 无关的群同态。若 \(D\) 是次数 \(d\ge0\) 的有效除子，则
	\begin{equation}\label{g4:eq:0402:degree-shift}
		A_q^{(d)}(D)=\operatorname{AJ}(D-d[q]).
	\end{equation}
	若整数一链 \(\Gamma\) 满足 \(\partial\Gamma=\sum_p n_p[p]\)，则同一个类也可写为
	\[
		\operatorname{AJ}\left(\sum_p n_p[p]\right)
		=\left[\left(\int_\Gamma\omega_1,\ldots,
		\int_\Gamma\omega_g\right)^{\mathsf T}\right].
	\]
\end{proposition}
\begin{proof}
	对每个支撑点选从 \(q\) 到它的道路，并乘以系数 \(n_p\) 后相加，得到一链 \(\Gamma\)。其边界为
	\[
		\partial\Gamma=\sum_p n_p([p]-[q])
		=\sum_p n_p[p],
	\]
	最后一步正使用次数为零。两条这样的链之差为整数一圈，积分向量之差属于周期格，所以商类不依赖链的选择。链相加对应积分相加，得到群同态。

	基点由 \(q\) 改为 \(q'\) 时，式\eqref{g4:eq:0402:basepoint-change}说明总变化为
	\[
		-\left(\sum_p n_p\right)A_q(q')=0.
	\]
	最后 \(A_q(q)=0\)，把 \(D\) 的各点逐项代入即得式\eqref{g4:eq:0402:degree-shift}。
\end{proof}

主除子都为零次除子，故有子群
\[
	\operatorname{Prin}(X)=
	\{\operatorname{div}(f):f\text{ 是 }X\text{ 上的非零亚纯函数}\}
	\subset\operatorname{Div}^0(X).
\]
除子积分提出一个具体问题：周期商类能否恰好检测两个同次数除子是否线性等价？下一节的 Abel 定理，即\cref{g4:thm:0403:abel}，将证明
\begin{equation}\label{g4:eq:0402:abel-kernel-preview}
	\ker\operatorname{AJ}=\operatorname{Prin}(X).
\end{equation}
该结论的两个方向都将在那里用第三类微分和指数积分证明。

\ProseParagraphBreak

先把这个结论的几何含义说清。令
\[
	\operatorname{Pic}^0(X)\coloneqq
	\operatorname{Div}^0(X)/\operatorname{Prin}(X).
\]
由\cref{g4:prop:0301:principal-trivial,g4:thm:0304:meromorphic-section}，它也就是零次全纯线丛的同构类群，运算为张量积：每条线丛可由非零亚纯截面的除子表示，而两个除子给出同构线丛当且仅当它们线性等价。一旦式\eqref{g4:eq:0402:abel-kernel-preview}证明，\(\operatorname{AJ}\) 就下降为单射
\[
	\operatorname{Pic}^0(X)\longrightarrow J(X).
\]
第4节的 Jacobi 反演将证明它满射，从而把零次除子类与复环面中的点对应起来。

\ProseParagraphBreak

对任意整数 \(d\)，记 \(\operatorname{Pic}^d(X)\) 为次数 \(d\) 的除子模线性等价所得的集合。选定 \(q\) 后，\(D\mapsto D-d[q]\) 把它识别为 \(\operatorname{Pic}^0(X)\) 的一个平移副本。换基点使对应目标坐标平移 \(-dA_q(q')\)。对于有效除子，这个几何解释已由实际构造的映射 \(A_q^{(d)}\) 承载；不同点组何时落到同一点，正是下一节要解决的线性等价问题。

\begin{example}
	\label{g4:exm:0402:sphere-divisors}
	球面亏格为零，故 \(J(\mathbb P^1)\) 是一点。若
	\[
		D=\sum_{a\in\mathbb C}n_a[a]+n_\infty[\infty],
		\qquad \sum_a n_a+n_\infty=0,
	\]
	其中只有有限项非零，则有理函数
	\[
		f(z)=\prod_{a\in\mathbb C}(z-a)^{n_a}
	\]
	满足 \(\operatorname{div}(f)=D\)。因为有限点的阶数为 \(n_a\)，无穷远阶数为 \(-\sum_a n_a=n_\infty\)。这直接说明球面上的零次除子都为主除子，与这里目标只有一点相符。
\end{example}
% 纲要条目排为编号小节标题，并进入目录与书签；说明段仍保留于纲要注释中。
% !TeX root = ../../Book4.tex
% 纲要标题：### 4.3 Abel 定理的证明
% 纲要位置：几何学计划纲要.md，第 3606 行。
% 正文状态：初稿；尚未正式审核。

% 纲要内容开始
% 
% * 主除子的必要周期条件
% * 第三类微分的构造
% * 整数周期与指数积分
% * 由周期条件恢复亚纯函数
% 
% 纲要内容结束

\section{Abel 定理的证明}
\label{g4:sec:0403}

给定度零除子 \(D=\sum_p n_p[p]\)，要构造满足 \(\operatorname{div}(f)=D\) 的亚纯函数，可以先寻找它的对数微分。若 \(f\) 在 \(p\) 的阶数为 \(n_p\)，则 \(\mathrm{d}f/f\) 在那里只留下留数为 \(n_p\) 的单极点。第三章的主部判据能构造具有这些留数的微分；真正的全局问题是它的积分能否取指数后成为单值函数。Abel 定理说明，这个问题恰由 \(D\) 在 Jacobi 簇中的积分像决定。

\subsection{主除子的必要周期条件}
\label{g4:subsec:0403:01}
先固定前两节的约定。设 \(X\) 为亏格 \(g\) 的紧连通黎曼曲面，取 \(a_i\cdot b_j=\delta_{ij}\) 的同调辛基及归一化全纯微分列向量
\[
	\boldsymbol\omega=(\omega_1,\ldots,\omega_g)^{\mathsf T},\qquad
	\int_{a_i}\omega_j=\delta_{ij},\qquad
	\tau_{ij}=\int_{b_i}\omega_j.
\]
当 \(g\ge1\) 时，由\cref{g4:thm:0401:normalized-periods,g4:prop:0401:period-lattice}，\(\tau=\tau^{\mathsf T}\)，且 \(\Lambda=\mathbb Z^g+\tau\mathbb Z^g\) 是周期格。\(g=0\) 时沿用 \(\mathbb C^0=\{0\}\)、\(\Lambda=\{0\}\) 的约定。度零除子的积分像记为
\[
	\operatorname{AJ}(D)
	=\left[\sum_p n_p\int_{q_0}^{p}\boldsymbol\omega\right]
	\in J(X)=\mathbb C^g/\Lambda.
\]
\cref{g4:prop:0402:divisor-integral}已经证明这个类与积分道路及基点 \(q_0\) 的选择无关；方括号内的复向量本身仍依赖道路。

\begin{proposition}
	\label{g4:prop:0403:logarithmic-differential}
	设 \(f\) 是紧连通黎曼曲面 \(X\) 上不恒为零的亚纯函数。其\term[duishuweifen]{对数微分}[logarithmic differential] \(\eta_f=\mathrm{d}f/f\) 在零极点之外全纯，在每个点至多有单极点，并满足
	\[
		\operatorname{Res}_p\eta_f=\operatorname{ord}_p(f).
	\]
	若闭路 \(\gamma\) 避开 \(f\) 的零极点，则
	\[
		\int_\gamma\eta_f\in2\pi\mathrm{i}\mathbb Z.
	\]
\end{proposition}
\begin{proof}
	在 \(p\) 的中心坐标中写 \(f=z^n u(z)\)，其中 \(n=\operatorname{ord}_p(f)\in\mathbb Z\)，且 \(u\) 全纯、\(u(0)\ne0\)。直接求微分得
	\[
		\frac{\mathrm{d}f}{f}
		=n\frac{\mathrm{d}z}{z}+\frac{u'(z)}{u(z)}\,\mathrm{d}z.
	\]
	第二项全纯，故得到局部极点与留数的结论。

	沿 \(\gamma\) 取有限个小邻域，在每个邻域选择非消失函数 \(f\) 的全纯对数，并在相邻段交接时加上 \(2\pi\mathrm{i}\) 的整数倍，使对数沿整条参数区间连续。各段上对数的微分都是 \(\mathrm{d}f/f\)，所以积分是终点对数值减去起点对数值。由于 \(f(\gamma(1))=f(\gamma(0))\)，这两个值的指数相同，其差属于 \(2\pi\mathrm{i}\mathbb Z\)。局部对数的存在使用\cref{g4:prop:0102:logarithm}在坐标圆盘上的结论。
\end{proof}

要把这些整数周期与除子的积分联系起来，需要在带极点的情形重新检查双线性公式。设有限点集 \(S\) 包含所考虑微分的全部极点。由\cref{g4:prop:0401:symplectic-basis}，可选几何辛基的环路代表，使它们避开 \(S\)，仍只在共同顶点相交，并沿这些环路把曲面切开成多边形 \(P\)。其正向边字为
\[
	a_1b_1a_1^{-1}b_1^{-1}\cdots
	a_gb_ga_g^{-1}b_g^{-1}.
\]
所有 \(p\in S\) 都位于 \(P\) 的内部。这里亚纯微分的 \(a,b\) 周期按这些实际环路取值：跨过极点改变环路时，周期可以改变一个小圆周积分；只说它们在 \(H_1(X,\mathbb Z)\) 中的类，还不足以确定亚纯微分的周期。

\ProseParagraphBreak

在 \(P\) 的内部、\(S\) 之外选基点 \(q_0\)。对全纯微分 \(\omega\)，令 \(F(p)=\int_{q_0}^{p}\omega\)，道路始终在切开后的多边形内。因为 \(P\) 是圆盘，\(F\) 是单值原函数，并在每条边的相应一侧延拓。特别地，它在各 \(p\in S\) 附近仍全纯。以下等式使用这一组确定的道路值。

\begin{proposition}
	\label{g4:prop:0403:reciprocity}
	设 \(X\) 为亏格 \(g\ge1\) 的紧连通黎曼曲面，\(\omega\) 为全纯微分，\(\eta\) 为至多在有限点集 \(S\) 有单极点的亚纯微分，记 \(\rho_p=\operatorname{Res}_p\eta\)。按上述方式选取避开 \(S\) 的切割多边形、其几何辛基及原函数 \(F\)。则
	\begin{equation}
		\label{g4:eq:0403:reciprocity}
		\sum_{i=1}^g
		\left(
			\int_{a_i}\omega\int_{b_i}\eta
			-\int_{b_i}\omega\int_{a_i}\eta
		\right)
		=2\pi\mathrm{i}\sum_{p\in S}\rho_pF(p).
	\end{equation}
	右边取实际复数值，不取商；改变 \(F\) 的加法常数不改变右边。
\end{proposition}
\begin{proof}
	在 \(P\) 内各 \(p\) 周围删去两两不交的小坐标圆盘，得到 \(P_\varepsilon\)。微分 \(F\eta\) 在这个带边界的区域全纯，且
	\[
		\mathrm{d}(F\eta)=\mathrm{d}F\wedge\eta+F\,\mathrm{d}\eta
		=\omega\wedge\eta=0.
	\]
	最后一个等式用到两条微分的局部表达都只有 \(\mathrm{d}z\) 项。Stokes 公式把外边界与各孔口联系起来。孔口作为 \(P_\varepsilon\) 的边界按顺时针定向，所以
	\[
		\int_{\partial P}F\eta
		=\sum_{p\in S}\int_{|z_p|=\varepsilon}^{\text{逆时针}}F\eta
		=2\pi\mathrm{i}\sum_{p\in S}\rho_pF(p).
	\]
	这里 \(\operatorname{Res}_p(F\eta)=F(p)\rho_p\)，因为 \(\eta\) 至多有单极点。多边形顶点处可先圆滑边角再取极限；\(F\) 与 \(\eta\) 在每个顶点副本附近有界，新增小弧上的积分趋于零，因此不产生额外项。

	再按配对边计算左边。简记 \(A_i=\int_{a_i}\omega\)、\(B_i=\int_{b_i}\omega\)。式\eqref{g4:eq:0401:edge-jumps}给出的原函数跳跃为：在同一 \(a_i\) 的两个匹配点上，负向边一侧的 \(F\) 比正向边一侧多 \(B_i\)；在同一 \(b_i\) 的两个匹配点上，正向边一侧的 \(F\) 比负向边一侧多 \(A_i\)。这些跳跃也可从一个四边块的边界行进直接读出。设块起点值为 \(C_i\)，沿 \(a_i,b_i\) 的部分积分分别为 \(f_a(t),f_b(t)\)，则匹配边上的值为
	\[
		\begin{array}{c|cc}
			&\text{正向边一侧}&\text{负向边一侧}\\ \hline
			a_i&C_i+f_a(t)&C_i+B_i+f_a(t)\\
			b_i&C_i+A_i+f_b(t)&C_i+f_b(t).
		\end{array}
	\]
	两侧的 \(\eta\) 来自同一条曲线上的同一个微分，而边界行进方向相反。于是这一对 \(a_i\) 边贡献 \(-B_i\int_{a_i}\eta\)，这一对 \(b_i\) 边贡献 \(A_i\int_{b_i}\eta\)。对全部边求和，便得到式\eqref{g4:eq:0403:reciprocity}的左边。

	最后，整体留数定理给出 \(\sum_p\rho_p=0\)。若 \(F\) 增加常数 \(C\)，右边只增加 \(2\pi\mathrm{i}C\sum_p\rho_p=0\)，所以不变。
\end{proof}

称式\eqref{g4:eq:0403:reciprocity}为这里的带极点互反公式。若所有留数都是整数 \(n_p\)，把 \(\omega\) 依次取为 \(\omega_1,\ldots,\omega_g\)，并记切盘内部道路所给出的列向量为
\[
	U(D)\coloneqq\sum_p n_p\int_{q_0}^{p}\boldsymbol\omega,
	\qquad D=\sum_p n_p[p],
\]
则互反公式写成
\begin{equation}
	\label{g4:eq:0403:vector-reciprocity}
	B(\eta)-\tau A(\eta)=2\pi\mathrm{i}U(D),
	\qquad
	A(\eta)=\left(\int_{a_i}\eta\right)_i,\quad
	B(\eta)=\left(\int_{b_i}\eta\right)_i.
\end{equation}
逐分量计算原本出现的是 \(\tau^{\mathsf T}\)，这里已使用其对称性。若重新任意选择到各点的道路，新的 \(U(D)\) 与原向量相差 \(\Lambda\) 中的元素；与固定环路周期之间的复向量等式则应继续使用切盘内确定的这一提升。

\begin{proposition}
	\label{g4:prop:0403:principal-periods}
	设 \(X\) 是紧连通黎曼曲面，\(f\) 是不恒为零的亚纯函数，则
	\[
		\operatorname{AJ}(\operatorname{div}(f))=0.
	\]
\end{proposition}
\begin{proof}
	由\cref{g4:prop:0203:principal-degree}，\(D=\operatorname{div}(f)\) 的次数为零。若 \(g=0\)，Jacobi 簇是零维的单点空间，结论成立。若 \(g\ge1\)，取 \(\eta=\mathrm{d}f/f\)，并选避开零极点的几何辛基。\cref{g4:prop:0403:logarithmic-differential}给出整数列向量 \(m,n\)，使
	\[
		A(\eta)=2\pi\mathrm{i}m,\qquad
		B(\eta)=2\pi\mathrm{i}n.
	\]
	式\eqref{g4:eq:0403:vector-reciprocity}于是给出
	\[
		U(D)=n-\tau m\in\mathbb Z^g+\tau\mathbb Z^g.
	\]
	它在商空间中的类正是 \(\operatorname{AJ}(D)\)，因此为零。若 \(f\) 为非零常数，则 \(D=0\)，亦有同样结论。
\end{proof}

\subsection{第三类微分的构造}
\label{g4:subsec:0403:02}
反方向先只处理极点附近的数据。设要在 \(p\) 取得留数 \(\rho\)，局部应取 \(\rho\,\mathrm{d}z/z\)。更换中心坐标 \(z=\zeta v(\zeta)\)、\(v(0)\ne0\)，则
\[
	\rho\frac{\mathrm{d}z}{z}
	=\rho\frac{\mathrm{d}\zeta}{\zeta}
	+\rho\frac{v'(\zeta)}{v(\zeta)}\,\mathrm{d}\zeta.
\]
后一项全纯，所以这一微分主部只依赖 \(p\) 与 \(\rho\)，不依赖中心坐标。只允许单极点的亚纯微分称为本节的\term[disanleiweifen]{第三类微分}[differential of the third kind]；若指定点的留数为零，则相应主部为零，微分在那里全纯。以下也允许全部主部为零的情形。

\begin{proposition}
	\label{g4:prop:0403:third-kind}
	设 \(X\) 是亏格 \(g\) 的紧连通黎曼曲面，\(S\subset X\) 有限，并指定复数 \(\rho_p\)、\(p\in S\)。存在只可能在 \(S\) 有单极点且留数为 \(\rho_p\) 的亚纯微分，当且仅当
	\[
		\sum_{p\in S}\rho_p=0.
	\]
	存在时，所有这样的微分组成 \(\eta_0+H^0(X,K_X)\)。若另外固定避开 \(S\) 的几何辛基，则恰有一条这样的微分 \(\eta_\rho\) 满足
	\[
		\int_{a_i}\eta_\rho=0\qquad(1\le i\le g).
	\]
	对由同一切割多边形确定的路径，归一化微分满足
	\begin{equation}
		\label{g4:eq:0403:third-kind-b-periods}
		\int_{b_j}\eta_\rho
		=2\pi\mathrm{i}\sum_{p\in S}\rho_p
		\int_{q_0}^{p}\omega_j
		\qquad(1\le j\le g).
	\end{equation}
\end{proposition}
\begin{proof}
	必要性是\cref{g4:thm:0203:global-residue}。对充分性，在每个 \(p\) 指定 \(K_X\) 值主部 \(\rho_p\,\mathrm{d}z_p/z_p\)，并应用\cref{g4:cor:0303:principal-parts-criterion}。这里取 \(L=K_X\)，对偶测试空间为
	\[
		H^0(X,K_X\otimes K_X^{-1})
		=H^0(X,\mathcal O_X)=\mathbb C.
	\]
	因此全部障碍条件就是：对每个常数 \(c\)，有
	\[
		\sum_{p\in S}\operatorname{Res}_p
		\left(c\rho_p\frac{\mathrm{d}z_p}{z_p}\right)
		=c\sum_{p\in S}\rho_p=0.
	\]
	所给条件恰好使它们全部成立，所以存在全局亚纯微分 \(\eta_0\)，具有指定主部，并在其余点全纯。因此总留数为零正是这一主部问题的全部障碍。

	两条这样的微分相减后，各单极点主部相消，差为全纯微分；反过来，加上全纯微分不改变主部。因此所有解确为所述仿射空间。令
	\[
		\eta_\rho=\eta_0-
		\sum_{i=1}^g\left(\int_{a_i}\eta_0\right)\omega_i.
	\]
	归一化条件 \(\int_{a_i}\omega_j=\delta_{ij}\) 使它的全部 \(a\) 周期为零。若另一条解也有零 \(a\) 周期，则差是一条全部 \(a\) 周期为零的全纯微分；由\cref{g4:thm:0401:normalized-periods}，这个差为零，故归一化解唯一。当 \(g=0\) 时，归一化条件为空，而 \(H^0(X,K_X)=0\)，同样得到唯一性。

	当 \(g\ge1\) 时，在\cref{g4:prop:0403:reciprocity}中依次取 \(\omega=\omega_j\)、\(\eta=\eta_\rho\)。全部 \(a\) 周期为零，左边只剩 \(\int_{b_j}\eta_\rho\)，从而得到式\eqref{g4:eq:0403:third-kind-b-periods}。\(g=0\) 时没有这样的周期等式。
\end{proof}

例如在 \(\mathbb P^1\) 上，给两个不同有限点 \(p,q\) 留数 \(1,-1\)，所得唯一微分为
\[
	\eta_{p,q}
	=\left(\frac1{z-p}-\frac1{z-q}\right)\mathrm{d}z
	=\frac{\mathrm{d}\bigl((z-p)/(z-q)\bigr)}{(z-p)/(z-q)}.
\]
它在无穷远全纯，因为两个单极点项的 \(z^{-1}\) 系数相消。若 \(q=\infty\)，则取 \(\mathrm{d}z/(z-p)\)，其无穷远留数为 \(-1\)。正亏格时，相同局部留数数据还允许加上 \(g\) 维的全纯微分；零 \(a\) 周期选定其中的一条，但它的 \(b\) 周期通常不在 \(2\pi\mathrm{i}\mathbb Z\) 中。

\ProseParagraphBreak

对度零整数除子 \(D=\sum_p n_p[p]\)，将这条归一化微分记为 \(\eta_D\)。式\eqref{g4:eq:0403:third-kind-b-periods}说明
\[
	A(\eta_D)=0,\qquad
	B(\eta_D)=2\pi\mathrm{i}U(D),
\]
其中 \(U(D)\) 是所选切割域中的实际积分向量。要使它成为某个函数的对数微分，还需要把这组 \(b\) 周期变为整数周期；允许的调整正是减去全纯微分。

\subsection{整数周期与指数积分}
\label{g4:subsec:0403:03}
在单连通区域，取原函数再取指数总能产生非消失全纯函数。去掉极点后的曲面通常不单连通；沿闭路继续原函数时会增加一个周期。指数函数恰好消去 \(2\pi\mathrm{i}\mathbb Z\) 中的增量，因此要检查的是全部闭路的周期。

\begin{lemma}
	\label{g4:lem:0403:punctured-periods}
	设 \(X\) 是亏格 \(g\) 的紧连通黎曼曲面，\(S=\{p_1,\ldots,p_r\}\) 为非空有限集。取避开 \(S\) 的几何辛基 \(a_i,b_i\)，并在每个 \(p_j\) 周围取逆时针小圆周 \(c_j\)。若 \(\xi\) 是 \(X\setminus S\) 上的闭复值光滑一形式，则任一闭路的 \(\xi\) 积分都是
	\[
		\int_{a_i}\xi,\qquad\int_{b_i}\xi,\qquad\int_{c_j}\xi
	\]
	的整数线性组合。
\end{lemma}
\begin{proof}
	先删去各点周围的小开圆盘。去心圆盘沿半径向外收缩到其边界，给出 \(X\setminus S\) 到所得带边界曲面的形变收缩，所以闭路可移到小圆盘之外。

	当 \(g\ge1\) 时，沿几何辛基切开，得到一个有 \(r\) 个孔口的多边形。在这个平面区域内，从各孔口向外边界依次选互不交叉的连接弧；每次切开一条弧，就把一个孔口接到外边界，孔的数目减少一个。全部切开后得到圆盘。当 \(g=0\) 时，底层曲面是球面，将一个孔口作为外边界，剩余区域是有 \(r-1\) 个孔的圆盘，按同样方式连接并切开即可；这时没有 \(a,b\) 环路。

	令 \(G\) 为切口图、各孔口圆周及连接弧在原带边界曲面中的并，细分弧的端点后它是有限连通图。收缩连接各部分的树状路径后，图的环路由 \(a_i,b_i\) 以及“走向孔口、绕孔一周、再返回”的环路生成。这是第一卷命题5.2.6的有限图计算：生成树之外每条边贡献一条生成环路。剩余的圆盘附着到 \(G\) 就恢复原曲面；第一卷命题5.2.7说明，附着圆盘只增加关系，图的基本群到所得空间基本群的同态满射。因此任意闭路的类均可用上述环路的有限乘积表示。

	\secref{g4:sec:0401}中闭形式积分的同调不变性论证也适用于无边界的 \(X\setminus S\)：每条光滑奇异链只有有限个单形，链上 Stokes 公式并不要求底空间紧。所以同伦替换保持积分，拼接时积分相加，反向时变号。连接孔口的路径在去程与回程上的积分相消，其贡献就是 \(\int_{c_j}\xi\)。这便给出所需整数线性组合。
\end{proof}

对具有整数留数 \(n_p\) 的第三类微分，小圆周积分已经自动等于 \(2\pi\mathrm{i}n_p\)。因此只需再把 \(a,b\) 周期调整到同一整数格中。没有极点时，前两节的同调辛基已生成全部周期，结论相同。

\begin{proposition}
	\label{g4:prop:0403:exponential-integral}
	设 \(X\) 是紧连通黎曼曲面，\(\eta\) 是至多有单极点、且所有留数均为整数的亚纯微分。则存在不恒为零的亚纯函数 \(f\)，使
	\[
		\frac{\mathrm{d}f}{f}=\eta,
	\]
	当且仅当 \(\eta\) 在其极点补集上沿每条闭路的积分均属于 \(2\pi\mathrm{i}\mathbb Z\)。这个条件等价于在一组避开极点的几何辛基上有
	\[
		\int_{a_i}\eta,\ \int_{b_i}\eta\in2\pi\mathrm{i}\mathbb Z
		\qquad(1\le i\le g).
	\]
	存在时，\(f\) 由 \(\eta\) 唯一确定到非零常数倍，且
	\[
		\operatorname{div}(f)
		=\sum_p\operatorname{Res}_p(\eta)[p].
	\]
\end{proposition}
\begin{proof}
	必要性由\cref{g4:prop:0403:logarithmic-differential}给出。小圆周积分由留数公式确定，故全部闭路与几何辛基上的条件等价，正是前一引理的结论。若没有极点，则用 \(H_1(X,\mathbb Z)\) 的辛基得到相同结论。

	现在设全部闭路满足整数周期条件，令 \(S\) 为极点集。在连通的 \(X\setminus S\) 中固定 \(q_0\)，定义
	\begin{equation}
		\label{g4:eq:0403:exponential-integral}
		f(p)\coloneqq
		\exp\left(\int_{q_0}^{p}\eta\right).
	\end{equation}
	删去有限点仍保持道路连通，因为经过某个删去点的道路可在小坐标盘内绕开它。两条从 \(q_0\) 到 \(p\) 的道路之差是闭路，积分差在 \(2\pi\mathrm{i}\mathbb Z\) 中，所以式\eqref{g4:eq:0403:exponential-integral}与道路选择无关。在不含极点的小盘中，\(\eta\) 有全纯原函数，因此 \(f\) 局部是某个全纯函数的指数乘以非零常数，故全纯、处处非零，并满足 \(\mathrm{d}f/f=\eta\)。

	只需检查它在 \(p\in S\) 的延拓。取中心坐标 \(z\)，写
	\[
		\eta=n_p\frac{\mathrm{d}z}{z}+h(z)\,\mathrm{d}z,
		\qquad n_p=\operatorname{Res}_p\eta\in\mathbb Z,
	\]
	并取全纯函数 \(H\)，使 \(H'=h\)。在去心盘上，
	\[
		\frac{\mathrm{d}}{\mathrm{d}z}
		\left(f(z)z^{-n_p}e^{-H(z)}\right)=0.
	\]
	去心盘连通，所以括号内等于某个非零常数 \(C_p\)。于是
	\begin{equation}
		\label{g4:eq:0403:local-exponential}
		f(z)=C_p z^{n_p}e^{H(z)}.
	\end{equation}
	整数幂在原点是单值亚纯的，而指数因子全纯且非零，因此 \(f\) 亚纯延拓，阶数恰为 \(n_p\)。这也给出所述除子。

	若 \(f_1,f_2\) 具有相同的对数微分，则它们的比值在极点补集上导数为零，因连通而为非零常数；等式再跨过有限点延拓。所以解唯一到非零常数倍。
\end{proof}

局部留数的整数性与全局周期条件各有作用。若留数为非整数 \(\rho\)，局部指数积分含 \(z^\rho\)，绕点一周会乘以 \(e^{2\pi\mathrm{i}\rho}\)。即使所有留数都为整数，沿把手环路继续积分仍可能产生非整数周期；这正是下一步要用全纯微分修正的部分。

\subsection{由周期条件恢复亚纯函数}
\label{g4:subsec:0403:04}
归一化第三类微分已经解决了全部局部主部条件。剩下的 \(g\) 个复周期恰由 \(U(D)\) 记录；当这个向量落在周期格中时，格的两组整数系数分别告诉我们应减去哪条全纯微分，以及修正后的另一组周期是多少。

\begin{theorem}[Abel 定理]
	\label{g4:thm:0403:abel}
	设 \(X\) 为紧连通黎曼曲面，\(D\) 是 \(X\) 上次数为零的整数除子。则下列两项等价：
	\begin{enumerate}
		\item 存在不恒为零的亚纯函数 \(f\)，使 \(\operatorname{div}(f)=D\)；
		\item \(\operatorname{AJ}(D)=0\in J(X)\)。
	\end{enumerate}
	在存在时，具有指定除子的函数唯一到非零常数倍。因此，对任意同次数除子 \(D_1,D_2\)，二者线性等价当且仅当 \(\operatorname{AJ}(D_1-D_2)=0\)。
\end{theorem}
\begin{proof}
	主除子的必要条件已由\cref{g4:prop:0403:principal-periods}证明。下面设 \(D=\sum_p n_p[p]\)、\(\sum_p n_p=0\)，并且 \(\operatorname{AJ}(D)=0\)。

	若 \(D=0\)，取 \(f=1\)。若 \(D\ne0\)，由\cref{g4:prop:0403:third-kind}构造只具有留数 \(n_p\) 的单极点、且全部 \(a\) 周期为零的微分 \(\eta_D\)。先设 \(g\ge1\)，在避开这些点的切割多边形内，选择前述确定的道路，得到 \(U(D)\in\mathbb C^g\)。其商类为零，意味着存在整数列向量 \(m,\ell\in\mathbb Z^g\)，使
	\[
		U(D)=m+\tau\ell.
	\]
	定义
	\begin{equation}
		\label{g4:eq:0403:integer-period-correction}
		\eta=\eta_D-2\pi\mathrm{i}
		\sum_{j=1}^g\ell_j\omega_j.
	\end{equation}
	所减微分全纯，所以 \(\eta\) 的极点及留数与 \(\eta_D\) 相同。其周期为
	\[
		A(\eta)=-2\pi\mathrm{i}\ell,\qquad
		B(\eta)
		=2\pi\mathrm{i}U(D)-2\pi\mathrm{i}\tau\ell
		=2\pi\mathrm{i}m.
	\]
	二者均在 \(2\pi\mathrm{i}\mathbb Z^g\) 中，孔口周期又为 \(2\pi\mathrm{i}n_p\)。于是\cref{g4:prop:0403:exponential-integral}适用，指数积分给出全局亚纯函数 \(f\)，并有 \(\operatorname{ord}_p(f)=n_p\)，所以 \(\operatorname{div}(f)=D\)。

	若 \(g=0\)，没有 \(a,b\) 周期，也没有需要减去的全纯微分。直接取 \(\eta=\eta_D\)；小圆周周期已为整数周期，\cref{g4:lem:0403:punctured-periods}使全部闭路满足条件，同样得到 \(f\)。这与 \(J(X)\) 为单点、每个度零除子的积分像都为零相符。

	若 \(f_1,f_2\) 具有相同除子，则 \(f_1/f_2\) 的所有零极点都相消，成为紧曲面上处处非零的全纯函数，故为非零常数。最后，\(D_1,D_2\) 线性等价按定义恰好表示 \(D_1-D_2\) 为主除子，再应用刚证得的等价关系即可。
\end{proof}

证明中的路径选择有两个层次。切割域提供一个确定的复向量 \(U(D)\)，使它与 \(\eta_D\) 的周期之间有精确等式；Abel--Jacobi 映射只保留其模 \(\Lambda\) 的类。零类保证这个特定提升也可写成 \(m+\tau\ell\)，从而允许上述周期修正。最终得到的函数虽然用到了这些选择，其除子只由 \(D\) 决定，而唯一性使不同选择所得函数只差非零常数。

\ProseParagraphBreak

Abel 定理使积分真正下降到除子类群，并给出单射
\[
	\operatorname{Div}^0(X)/\operatorname{Prin}(X)
	\hookrightarrow J(X).
\]
它还可以用于检验一条曲线本身在 Jacobi 簇中的积分像。

\begin{corollary}
	\label{g4:cor:0403:abel-embedding}
	设 \(X\) 为亏格 \(g\ge1\) 的紧连通黎曼曲面，固定 \(q_0\in X\)。则
	\[
		A_{q_0}:X\longrightarrow J(X),\qquad
		p\longmapsto\left[\int_{q_0}^{p}\boldsymbol\omega\right]
	\]
	是全纯嵌入。若 \(g=1\)，它是将 \(q_0\) 送到原点的双全纯同构
	\[
		X\cong\mathbb C/(\mathbb Z+\tau\mathbb Z).
	\]
\end{corollary}
\begin{proof}
	若 \(A_{q_0}(p)=A_{q_0}(q)\)，则 \(\operatorname{AJ}([p]-[q])=0\)。Abel 定理给出具有除子 \([p]-[q]\) 的函数。若 \(p\ne q\)，这个函数恰有一个简单极点，到球面的次数为 \(1\)，由\cref{g4:cor:0202:degree-one}是双全纯同构。这会使 \(g=0\)，与假设矛盾。因此 \(p=q\)，映射单射。

	由\cref{g4:prop:0402:abel-jacobi}，其微分为
	\[
		\mathrm{d}(A_{q_0})_p(v)
		=(\omega_1(v),\ldots,\omega_g(v))^{\mathsf T}.
	\]
	当 \(g\ge2\) 时，\cref{g4:thm:0305:canonical-map}的典范系无基点性保证每个 \(p\) 至少有一条全纯微分不消失，所以这个微分在一维切空间上单射。当 \(g=1\) 时，取任一非零全纯微分；它的有效零除子次数为 \(\deg K_X=0\)，故处处非零，得到同样结论。

	于是 \(A_{q_0}\) 是单射全纯浸入。紧性使它成为到其像的同胚；局部取一个导数非零的目标坐标，一元逆函数定理使其余目标坐标成为这个坐标的全纯函数，所以像是复子流形，映射为全纯嵌入。若 \(g=1\)，源、靶维数同为 \(1\)，映射局部双全纯，像开；紧性又使像闭。\(J(X)\) 连通，因此像等于整个 Jacobi 簇。基点的积分为零，故 \(q_0\) 对应原点。
\end{proof}

在球面上，结论可由有理函数直接核验。若 \(D=\sum_{j=1}^r n_j[a_j]+n_\infty[\infty]\)，且 \(n_\infty=-\sum_jn_j\)，则
\[
	f(z)=\prod_{j=1}^r(z-a_j)^{n_j}
\]
恰有除子 \(D\)。在亏格一时，将曲面写成上述复环面，度零除子 \(\sum_j n_j[p_j]\) 为主除子当且仅当代表点满足
\[
	\sum_j n_j z_j\in\mathbb Z+\tau\mathbb Z.
\]
例如任取环面上的 \(u,v\)，除子 \([u]+[v]-[0]-[u+v]\) 的积分像为零，因而可由一个亚纯函数实现。周期格在这里直接规定了零点与极点在环面群法中的平衡条件。
% 纲要条目排为编号小节标题，并进入目录与书签；说明段仍保留于纲要注释中。

% !TeX root = ../../Book4.tex
% 纲要标题：### 4.4 Jacobi 反演及 theta 几何
% 纲要位置：几何学计划纲要.md，第 3613 行。
% 正文状态：初稿；尚未正式审核。

% 纲要内容开始
% 
% * 反演的满射性与 Riemann–Roch
% * 特殊除子及纤维维数
% * theta 除子和主极化
% * Torelli 定理与曲面重构
% 
% 纲要内容结束

\section{Jacobi 反演及 theta 几何}
\label{g4:sec:0404}

Abel 定理解决了积分映射何时把两个点组识别起来。现在反过来：给定 Jacobi 环面中的一个点，能否找到一个点组实现它？如果可以，需要多少个点，所有实现是否唯一？Riemann--Roch 把这些问题与线性系联系起来。积分矩阵还给出一条全纯级数，其零点形成 theta 除子；它同时记录某些除子的有效性和周期格上的相交配对。

\subsection{反演的满射性与 Riemann–Roch}
\label{g4:subsec:0404:01}

固定紧连通黎曼曲面 \(X\) 和基点 \(q\)，记亏格为 \(g\)。\(\operatorname{Sym}^dX\) 的点是次数为 \(d\) 的有效除子，\cref{g4:thm:0402:symmetric-product}已证明它是紧连通复 \(d\) 维流形。上一节的 Abel 定理立即给出
\begin{equation}
	\label{g4:eq:0404:same-fiber}
	A_q^{(d)}(D)=A_q^{(d)}(D')
	\quad\Longleftrightarrow\quad D\sim D'.
\end{equation}
这里两边除子必须有相同次数。一个纤维因而就是一个线性等价类中全部有效除子组成的集合。

\begin{lemma}
	\label{g4:lem:0404:nonspecial-divisor}
	设 \(X\) 为亏格 \(g\ge1\) 的紧连通黎曼曲面。存在互不相同的点 \(p_1,\ldots,p_g\)，使 \(D_0=[p_1]+\cdots+[p_g]\) 满足
	\[
		H^0(X,K_X(-D_0))=0,\qquad h^0(\mathcal O_X(D_0))=1.
	\]
	在 \(D_0\) 处，\(A_q^{(g)}\) 的微分是复线性同构，且其整个纤维只有 \(D_0\) 一个点。
\end{lemma}
\begin{proof}
	先从 \(g\) 维空间 \(V_0=H^0(X,K_X)\) 开始。只要 \(V_{r-1}\ne0\)，从中取一条非零微分；它在紧曲面上只有有限个零点，所以可以在此前各点之外选 \(p_r\)，使其在 \(p_r\) 非零。求值映射 \(V_{r-1}\to(K_X)_{p_r}\) 非零，目标一维，因而其核
	\[
		V_r=H^0\bigl(X,K_X(-p_1-\cdots-p_r)\bigr)
	\]
	恰减少一维。经过 \(g\) 步得到 \(V_g=0\)。Riemann--Roch 给出 \(h^0(\mathcal O(D_0))=g+1-g=1\)。

	由\cref{g4:prop:0402:abel-differential}，\(\mathrm{d}(A_q^{(g)})_{D_0}\) 的秩为 \(g-h^0(K_X(-D_0))=g\)，所以它是同构。若 \(D'\) 与 \(D_0\) 具有相同积分像，Abel 定理给出 \(D'=D_0+\operatorname{div}(f)\)，其中 \(0\ne f\in H^0(\mathcal O(D_0))\)。这个一维空间由常数张成，所以 \(D'=D_0\)。
\end{proof}

局部可逆只保证像中含一个开集；要得到全部目标，还要利用紧性与映射度。第二卷第16.1节定理16.1.4说明：同维连通闭定向流形之间的光滑映射，其同调度等于任意正则纤维中的局部定向符号之和。遗漏一个目标点时，那个点是空纤维的正则值，因而映射度为零。下面将它用于复流形的自然定向；复线性同构作为实线性映射的行列式为其复行列式模的平方，所以符号总为正。

\begin{theorem}[Jacobi 反演]
	\label{g4:thm:0404:jacobi-inversion}
	设 \(X\) 为亏格 \(g\) 的紧连通黎曼曲面，\(q\in X\)。映射
	\[
		A_q^{(g)}:\operatorname{Sym}^gX\longrightarrow J(X)
	\]
	满射。若 \(g\ge1\)，它作为实 \(2g\) 维闭定向流形之间的映射，映射度为 \(1\)。因此，任意 \(u\in J(X)\) 都能写成
	\[
		u=A_q(p_1)+\cdots+A_q(p_g),
	\]
	其中允许某些 \(p_i\) 重合。
\end{theorem}
\begin{proof}
	当 \(g=0\) 时，\(\operatorname{Sym}^0X\) 与 \(J(X)\) 都是一点，结论成立。设 \(g\ge1\)。源与目标均为紧连通复 \(g\) 维流形，故满足上述映射度定理的条件。

	取\cref{g4:lem:0404:nonspecial-divisor}中的 \(D_0\)，令 \(u_0=A_q^{(g)}(D_0)\)。这个值只有一个原像，且微分为复线性同构。因此 \(u_0\) 是正则值，其纤维符号和为 \(1\)，从而 \(\deg A_q^{(g)}=1\)。非零映射度排除了遗漏目标点的可能，所以映射满射。
\end{proof}

反演定理同时完成了除子类的参数化。任意 \(u\in J(X)\) 都是某个 \(\operatorname{AJ}(D-g[q])\)，其中 \(D\) 有效且次数为 \(g\)；因此 Abel 定理给出的单射实际上是群同构
\begin{equation}
	\label{g4:eq:0404:picard-jacobian}
	\operatorname{Pic}^0(X)\xrightarrow{\ \sim\ }J(X).
\end{equation}
对任意整数 \(d\)，先作 \(L\mapsto L(-d[q])\)，就得到 \(\operatorname{Pic}^d(X)\) 与同一个复环面的对应。这给 \(\operatorname{Pic}^d(X)\) 赋予复流形结构；更换 \(q\) 只作平移。\(\operatorname{Pic}^0\) 的零点是平凡线丛，其余次数的副本则需要选择基准线丛才有指定原点。

\subsection{特殊除子及纤维维数}
\label{g4:subsec:0404:02}

反演保证存在，却允许一个积分值由许多点组实现。这份不唯一性可以准确计数：固定一个有效除子 \(D\)，同一纤维中的另一个除子恰是 \(D+\operatorname{div}(f)\)，而可用函数 \(f\) 构成 \(H^0(X,\mathcal O(D))\)。非零常数倍不改变零点除子，所以应取这个向量空间的射影化。

\begin{theorem}
	\label{g4:thm:0404:abel-fibers}
	设 \(X\) 是亏格 \(g\) 的紧连通黎曼曲面，\(D\) 为次数 \(d\ge0\) 的有效除子。则 \(A_q^{(d)}\) 在 \(A_q^{(d)}(D)\) 上的纤维为完全线性系
	\[
		|D|\cong\mathbb P H^0(X,\mathcal O(D)).
	\]
	这个对应给出 \(|D|\) 到 \(\operatorname{Sym}^dX\) 的全纯嵌入。特别地，纤维的复维数为
	\begin{equation}
		\label{g4:eq:0404:fiber-dimension}
		\dim|D|=h^0(\mathcal O(D))-1
		=d-g+h^0(K_X(-D)).
	\end{equation}
	而 Abel 映射在 \(D\) 处的微分秩为 \(g-h^0(K_X(-D))\)。
\end{theorem}
\begin{proof}
	式\eqref{g4:eq:0404:same-fiber}已确定纤维的集合。把非零截面 \(s\) 送到其有效零除子 \(\operatorname{div}_{\mathcal O(D)}(s)\)，便得到映射
	\[
		\mathbb P H^0(X,\mathcal O(D))\longrightarrow\operatorname{Sym}^dX.
	\]
	它的像正是所求纤维。若两截面有同一个零除子，其比值是处处非零的全纯函数，紧性使它为非零常数，所以映射单射。

	说明全纯性时，在某个截面的各零点附近取互不相交的圆盘，边界避开零点。充分小的截面变动不会改变各盘内的零点总数，这是 Rouché 定理的应用。以 \(f_s(z)\) 表示截面的局部系数，各盘内零点的幂和由
	\[
		\frac1{2\pi\mathrm{i}}\int_{\partial U}
		z^k\frac{f_s'(z)}{f_s(z)}\,\mathrm{d}z
	\]
	给出，它全纯依赖截面参数。初等对称多项式由幂和递推得到，因此\cref{g4:thm:0402:symmetric-product}的坐标中，这个零除子映射全纯。

	还要检查浸入性。一个射影切向量由截面变化 \(s+t\dot s\) 模去 \(\mathbb C s\) 表示。若零除子的一阶变化为零，则每个零点盘内的所有初等对称坐标以及幂和的一阶变化均为零。在 \(s\) 有 \(m\) 阶零点的中心坐标中，令 \(r=\dot s/s\)；它的极点阶数至多为 \(m\)。对上式求导并沿圆周分部积分，得到对 \(1\le k\le m\)
	\[
		0=\frac1{2\pi\mathrm{i}}\int_{\partial U}z^k r'(z)\,\mathrm{d}z
		=-k\operatorname{Res}_0(z^{k-1}r(z)\,\mathrm{d}z).
	\]
	这些等式逐一消去 \(r\) 的全部主部系数。因此 \(\dot s/s\) 在所有零点处可去，并在整个紧曲面上全纯，故为常数。原射影切向量遂为零。映射是单射全纯浸入，源又紧，所以它是全纯嵌入。

	维数公式由 Riemann--Roch 得到；最后的微分秩已在\cref{g4:prop:0402:abel-differential}证明。两式相加为 \(d\)，也核对了纤维切空间的维数。
\end{proof}

称 \(D\) 特殊，仍是指 \(h^0(K_X(-D))>0\)。对 \(d=g\)，非特殊除子所在的纤维为一个点，且积分映射在该点局部双全纯；特殊除子所在的纤维则为正维射影空间。对 \(d\ge\max(0,2g-1)\)，次数使 \(K_X(-D)\) 无非零截面，所以所有非空纤维都是 \(\mathbb P^{d-g}\)。当 \(g=0,d=0\) 时同样得到 \(\mathbb P^0\)。这些高次数映射实际上满射：对任意次数 \(d\) 的线丛，Riemann--Roch 给出 \(h^0=d+1-g>0\)，其非零截面的零除子便实现相应的积分值。

\begin{example}
	\label{g4:exm:0404:genus-two-fiber}
	设 \(g=2\)。对次数 \(2\) 的有效除子 \(D\)，线丛 \(K_X(-D)\) 的次数为零。它有非零截面当且仅当平凡，即当且仅当 \(D\sim K_X\)。因此 \(A_q^{(2)}\) 的全部纤维中只有典范类对应一个 \(\mathbb P^1\)，其余纤维均为一点。这个 \(\mathbb P^1\) 正是超椭圆二重映射的全部纤维除子组成的线性系。它把反演不唯一的原因呈现为一个具体的特殊除子类。
\end{example}

\subsection{theta 除子和主极化}
\label{g4:subsec:0404:03}

以下设 \(g\ge1\)，写归一化周期矩阵为 \(\tau=U+\mathrm{i}Y\)，其中 \(\tau^{\mathsf T}=\tau\)、\(Y>0\)。\cref{g4:prop:0401:period-lattice}给出 \(J(X)=\mathbb C^g/(\mathbb Z^g+\tau\mathbb Z^g)\)。为了在这个商空间上构造除子，可以寻找一个全纯函数：它平移一个周期后只乘以一个处处非零的因子。函数本身未必下降，零点却能够下降。

\begin{definition}
	\label{g4:def:0404:theta}
	对称矩阵 \(\tau\in M_g(\mathbb C)\) 满足 \(\operatorname{Im}\tau>0\) 时，定义 \term[thetahanshu]{theta 函数}[theta function]
	\begin{equation}
		\label{g4:eq:0404:theta-series}
		\theta(z,\tau)=\sum_{n\in\mathbb Z^g}
		\exp\bigl(\pi\mathrm{i}n^{\mathsf T}\tau n+2\pi\mathrm{i}n^{\mathsf T}z\bigr),
		\qquad z\in\mathbb C^g.
	\end{equation}
\end{definition}

\begin{proposition}
	\label{g4:prop:0404:theta-automorphy}
	级数\eqref{g4:eq:0404:theta-series}及其任意阶 \(z\) 导数在紧集上一致绝对收敛，定义一个不恒为零的整函数。对 \(m,n\in\mathbb Z^g\)，有
	\begin{equation}
		\label{g4:eq:0404:theta-automorphy}
		\theta(z+m+\tau n,\tau)
		=\exp\bigl(-\pi\mathrm{i}n^{\mathsf T}\tau n-2\pi\mathrm{i}n^{\mathsf T}z\bigr)\theta(z,\tau).
	\end{equation}
\end{proposition}
\begin{proof}
	取 \(c>0\)，使 \(n^{\mathsf T}Yn\ge c|n|^2\)。当 \(z\) 位于固定紧集时，存在 \(M\) 使各项的绝对值不超过
	\[
		\exp(-\pi c|n|^2+2\pi M|n|).
	\]
	对 \(z\) 求任意阶导数，只多出一个关于 \(n\) 的多项式因子。高斯衰减使这些级数均可求和，因而可以逐项求导。这里多复变量的全纯性也可由逐坐标 Cauchy 公式及紧集一致收敛得到。

	整数平移 \(m\) 给各项乘以 \(e^{2\pi\mathrm{i}n^{\mathsf T}m}=1\)。对平移 \(\tau n\)，把求和指标由 \(k\) 改为 \(k+n\)，利用 \(\tau\) 对称完成平方，就得到式\eqref{g4:eq:0404:theta-automorphy}。在实单位立方体上积分，所有非零 Fourier 项都消失，而零项等于 \(1\)，所以
	\[
		\int_{[0,1]^g}\theta(s,\tau)\,\mathrm{d}s=1.
	\]
	因此 \(\theta\) 不恒为零。
\end{proof}

令 \(\lambda=m+\tau n\)，把式\eqref{g4:eq:0404:theta-automorphy}中的因子记为 \(j_\lambda(z)\)。它满足
\[
	j_{\lambda+\mu}(z)=j_\mu(z+\lambda)j_\lambda(z).
\]
确实，若 \(\mu=m'+\tau n'\)，两边指数的差是 \(-2\pi\mathrm{i}(n')^{\mathsf T}m\)，其指数为 \(1\)。于是关系
\[
	(z,v)\sim(z+\lambda,j_\lambda(z)v)
\]
把 \(\mathbb C^g\times\mathbb C\) 粘成 \(J(X)\) 上的全纯线丛 \(L_\theta\)。局部取不遇到非平凡格平移的小球，就得到其线丛图册。\(\theta\) 的变换律正说明它是 \(L_\theta\) 的全局全纯截面。它的局部零点方程在换图时相差一个全纯单位，因而定义有效的\term[thetachuzi]{theta 除子}[theta divisor]，记为 \(\Theta\)。

\begin{proposition}
	\label{g4:prop:0404:theta-section-dimension}
	上述线丛满足 \(h^0(J(X),L_\theta)=1\)，其截面空间由 \(\theta\) 张成。
\end{proposition}
\begin{proof}
	任一截面提升为整函数 \(f(z)\)，满足与 \(\theta\) 相同的变换律。特别地，\(f(z+m)=f(z)\)。逐坐标令 \(w_j=e^{2\pi\mathrm{i}z_j}\)，它下降为 \((\mathbb C^*)^g\) 上的全纯函数。对每个变量应用 Laurent 展开，再用 Cauchy 系数公式在紧多环域上控制级数，得到局部一致收敛的展开
	\[
		f(z)=\sum_{k\in\mathbb Z^g}c_k e^{2\pi\mathrm{i}k^{\mathsf T}z}.
	\]
	这一步只是把一元 Laurent 展开依次用于各变量；系数由各圆周上的重复积分唯一确定。

	将 \(z\mapsto z+\tau n\) 的变换律代入并比较 Fourier 系数，得到
	\[
		c_{k+n}=\exp\bigl(\pi\mathrm{i}n^{\mathsf T}\tau n
		+2\pi\mathrm{i}k^{\mathsf T}\tau n\bigr)c_k.
	\]
	取 \(k=0\)，所有系数均由 \(c_0\) 决定，且恰是 \(c_0\theta\) 的系数。因此截面空间至多一维；\(\theta\ne0\) 又保证它恰为一维。
\end{proof}

\begin{definition}
	\label{g4:def:0404:principal-polarization}
	设 \(V/\Lambda\) 为复环面。若格上的整数交替双线性形式 \(E\) 延伸到实向量空间 \(V\) 后满足
	\[
		E(\mathrm{i}v,\mathrm{i}w)=E(v,w),\qquad
		E(v,\mathrm{i}v)>0\quad(v\ne0),
	\]
	则称 \(E\) 为这个复环面的一个\term[jihua]{极化}[polarization]。若
	\[
		\Lambda\longrightarrow\operatorname{Hom}(\Lambda,\mathbb Z),\qquad
		\lambda\longmapsto E(\lambda,\cdot)
	\]
	是同构，则称为\term[zhujihua]{主极化}[principal polarization]。在任一整数格基中，后一个条件等价于 \(E\) 的矩阵行列式为 \(1\)。
\end{definition}

这里采用复环面的解析表述：整数性使正性资料能够由线丛承载。Jacobi 环面的相交形式提供一个标准主极化，下面把它与刚构造的 theta 线丛直接对应。用 \(y=\operatorname{Im}z\)，在覆盖空间上设
\[
	h(z)=\exp\bigl(-2\pi y^{\mathsf T}Y^{-1}y\bigr).
\]
它满足
\[
	h(z+\lambda)|j_\lambda(z)|^2=h(z),
\]
因为两边指数中有关 \(n^{\mathsf T}Yn\) 和 \(n^{\mathsf T}y\) 的项相消。因此 \(|v|^2h(z)\) 给 \(L_\theta\) 一个光滑 Hermitian 度量。

\begin{theorem}
	\label{g4:thm:0404:theta-polarization}
	Jacobi 环面的标准交替形式
	\[
		E(v,w)=-\operatorname{Im}\bigl(v^{\mathsf T}Y^{-1}\overline w\bigr)
	\]
	是主极化。theta 线丛的一阶陈类\footnote{陈省身（Shiing-Shen Chern，1911--2004），华裔美籍数学家，研究微分几何与示性类，创立陈类理论。}的实像由正闭二形式
	\begin{equation}
		\label{g4:eq:0404:theta-curvature}
		\Omega=\frac{\mathrm{i}}2\sum_{j,k}(Y^{-1})_{jk}
		\,\mathrm{d}z_j\wedge\mathrm{d}\overline z_k
	\end{equation}
	代表，且 \(\Omega(v,w)=E(v,w)\)。特别地，\(\Theta\) 非空，并有
	\[
		\int_{J(X)}\frac{\Omega^g}{g!}=1.
	\]
\end{theorem}
\begin{proof}
	由\eqref{g4:eq:0401:polarization-form}，\(E\) 与复结构相容，且 \(E(v,\mathrm{i}v)=v^{\mathsf T}Y^{-1}\overline v>0\)。由\eqref{g4:eq:0401:integral-alternating-form}，在格基 \((e_1,\ldots,e_g,\tau e_1,\ldots,\tau e_g)\) 下，它的矩阵是 \(\left(\begin{smallmatrix}0&I\\-I&0\end{smallmatrix}\right)\)，所以整数性与幺模性同时成立。

	为核对线丛，在局部全纯标架中用度量系数 \(h\) 定义联络形式 \(A=\partial\log h\)。若标架变为 \(e'=a e\)，则 \(h'=|a|^2h\)，于是 \(A'=A+a^{-1}\mathrm{d}a\)；这正是联络的变换律。等式 \(A+\overline A=\mathrm{d}\log h\) 验证度量相容，所以这是酉联络。其曲率为 \(F=\bar\partial\partial\log h=-\partial\bar\partial\log h\)。第二卷第22.3节命题22.3.1证明 \(\mathrm{i}F/(2\pi)\) 代表 \(c_1(L_\theta)_{\mathbb R}\)。对当前二次函数直接求导，得到
	\[
		\frac{\mathrm{i}F}{2\pi}
		=-\frac{\mathrm{i}}{2\pi}\partial\bar\partial\log h
		=\frac{\mathrm{i}}2\sum_{j,k}(Y^{-1})_{jk}
		\,\mathrm{d}z_j\wedge\mathrm{d}\overline z_k=\Omega.
	\]
	在实向量 \(v,w\) 上代入右式，得到 \(-\operatorname{Im}(v^{\mathsf T}Y^{-1}\overline w)\)，所以它正是 \(E\) 对应的常系数二形式。

	在实坐标 \(z=s+\tau t\) 中，\(\Omega=\sum_j\mathrm{d}s_j\wedge\mathrm{d}t_j\)。它在由任意一对 \(e_j,\tau e_j\) 张成的二维子环面上积分为 \(1\)，从而其上同调类非零。若 \(\theta\) 没有零点，便给出 \(L_\theta\) 的全局全纯标架，使该线丛平凡、一阶陈类为零，矛盾。因此 \(\Theta\) 非空。最后在基本平行多面体上按复定向积分 \(\Omega^g/g!\)，得到 \(g\) 个单位二维面积元的乘积，值为 \(1\)。
\end{proof}

线丛由乘子表达，主极化则只保留其整数二次上同调类。两者的关系还可以从除子观察：\(h^0(L_\theta)=1\) 说明固定这条线丛后，非零截面的零除子唯一。对线丛及其截面同时作平移，会使零除子平移，而标准极化类不变；同一线丛仅更换平凡化时，零除子保持不变。

\ProseParagraphBreak

有效除子也在 Jacobi 环面中确定一列紧子集
\[
	W_d=A_q^{(d)}(\operatorname{Sym}^dX),\qquad d\ge0.
\]
Jacobi 反演给出 \(W_g=J(X)\)；加上基点还给出 \(W_d=J(X)\) 对所有 \(d\ge g\) 成立。恰好少一个点的 \(W_{g-1}\) 则是 theta 几何与线性系相接的位置。下述更强的除子识别采用文献中的定理，其证明见\cite[第2章第7节，第338--340页，Riemann's Theorem]{GriffithsHarris1994}；前置的 theta 截面拉回与常数计算在同书第336--337页。

\begin{theorem}[Riemann 的 theta 定理]
	\label{g4:thm:0404:riemann-theta}
	设 \(X\) 为亏格 \(g\ge1\) 的紧连通黎曼曲面，固定上述归一化周期与基点 \(q\)。存在常向量类 \(\kappa_q\in J(X)\)，使作为有效除子有
	\[
		\Theta=W_{g-1}+\kappa_q.
	\]
	右边的 \(W_{g-1}\) 取其约化有效除子结构。等价地，\(u\in J(X)\) 满足 \(\theta(u+\kappa_q,\tau)=0\)，当且仅当次数 \(g-1\) 的相应线丛有非零全纯截面。
\end{theorem}

式中的 theta 等式按格商解释，其零性不依赖代表。这个定理把级数零点与除子的有效性相认；本节此前的反演、纤维维数及主极化计算各有独立证明。对于 \(g=1\)，\(W_0=\{0\}\)，所以 theta 除子是一个平移后的点；对于 \(g=2\)，它是 Abel 嵌入曲线的一个平移。更高亏格时，它把所有可由 \(g-1\) 个点表示的积分值组织成一个除子。

\subsection{Torelli 定理与曲面重构}
\label{g4:subsec:0404:04}

本章至此从曲面构造了复环面 \(J(X)\)、整数主极化 \(E_X\)，以及由基点选择决定平移的 theta 除子。反演是固定一条曲面之后在这个环面内寻找点组；重构问题则问这些整体资料是否足以找回原曲面。Torelli 问题\footnote{Ruggiero Torelli（1884--1915），意大利数学家，研究代数曲线与 Jacobi 簇。}要求保留主极化：给定两个曲面，比较的是 \((J(X),E_X)\) 与 \((J(Y),E_Y)\)。

\ProseParagraphBreak

具体地，一个保持原点的复环面同构可由复线性同构 \(C:\mathbb C^g\to\mathbb C^g\) 给出，满足 \(C\Lambda_X=\Lambda_Y\)。它保持主极化，是指
\[
	E_Y(C\lambda,C\mu)=E_X(\lambda,\mu)
	\qquad(\lambda,\mu\in\Lambda_X).
\]
若从一个保持原点的全纯群同构 \(F\) 出发，也能得到这样的 \(C\)。把两个环面的原点切空间识别为各自的覆盖向量空间，令 \(C=\mathrm{d}F_0\)。群同态条件说明 \(\mathrm{d}F\) 在平移识别下处处等于 \(C\)。所以从 \(\mathbb C^g\) 到目标环面的映射 \(z\mapsto F([z])-[Cz]\) 导数处处为零，又在原点取零，故恒为零。于是 \(F([z])=[Cz]\)，并有 \(C\Lambda_X\subseteq\Lambda_Y\)；再对逆映射应用同样结论，便得格的相等。

\begin{theorem}[经典 Torelli 定理]
	\label{g4:thm:0404:torelli}
	设 \(X,Y\) 是紧连通黎曼曲面，亏格至少为 \(1\)。若它们的 Jacobi 簇存在保持原点及标准主极化的全纯群同构，则 \(X,Y\) 双全纯同构。反过来，曲面的双全纯同构诱导其主极化 Jacobi 簇同构。
\end{theorem}

这个重构定理在此引用，证明见\cite[第2章第7节，第359--362页，Torelli's Theorem]{GriffithsHarris1994}。反方向可以从本章构造直接看出：双全纯映射拉回全纯微分、推前整数同调，积分与二者相容；它又保持复定向及曲线相交数，所以保持标准极化。正方向需要进一步分析 theta 除子和典范曲线的关系，特别要处理超椭圆情形。亏格 \(1\) 时，\cref{g4:cor:0403:abel-embedding}已把曲面本身与其 Jacobi 环面识别；下一节将把这个识别写成椭圆函数与三次方程。
% 纲要条目排为编号小节标题，并进入目录与书签；说明段仍保留于纲要注释中。

% !TeX root = ../../Book4.tex
% 纲要标题：### 4.5 椭圆曲线的解析与代数模型
% 纲要位置：几何学计划纲要.md，第 3620 行。
% 正文状态：初稿；尚未正式审核。

% 纲要内容开始
% 
% * Weierstrass 函数及微分方程
% * 复环面到光滑三次曲线
% * 加法公式与几何群律
% * 格变换、j 不变量与退化
% 
% ---
% 
% 纲要内容结束

\section{椭圆曲线的解析与代数模型}
\label{g4:sec:0405}

第三章从极点阶数为 \(2,3\) 的函数构造了亏格一曲面的三次模型。本章的周期积分又把同一曲面写成 \(\mathbb C/\Lambda\)。现在要把这两种描述直接接起来：从格上的一个级数出发，找出那两个函数，写出它们满足的方程，并把环面上的加法译成三次曲线上的弦切作图。最后，改变格与改变方程系数将给出同一个复参数 \(j\)。

\subsection{Weierstrass 函数及微分方程}
\label{g4:subsec:0405:01}

设
\[
	\Lambda=\mathbb Z\lambda_1+\mathbb Z\lambda_2\subset\mathbb C,
	\qquad \operatorname{Im}(\lambda_2/\lambda_1)>0.
\]
这里 \(\lambda_1,\lambda_2\) 是完整的周期，而非半周期。商 \(T_\Lambda=\mathbb C/\Lambda\) 由\cref{g4:prop:0402:complex-torus}成为紧连通黎曼曲面。一个以 \(\Lambda\) 为周期的亚纯函数称为关于该格的\term[tuoyuanhanshu]{椭圆函数}[elliptic function]；它等价于 \(T_\Lambda\) 上的亚纯函数。整的椭圆函数必为常数，因为它在一个紧基本平行四边形上有界，周期性使此界扩展到整个复平面，再由\cref{g4:prop:0101:cauchy-estimate}得 \(|f'(z)|\le M/r\)，令 \(r\to\infty\) 即得 \(f'=0\)。

\ProseParagraphBreak

希望在每个格点放一个相同的二阶极点，最直接的和是 \(\sum_{\lambda\in\Lambda}(z-\lambda)^{-2}\)，但其绝对收敛性不足。减去每项在 \(z=0\) 的常数值，便能多获得一阶衰减，并同时规定 Laurent 展开的常数项为零。

\begin{theorem}\label{g4:thm:0405:weierstrass-series}
	设 \(\Lambda\subset\mathbb C\) 为秩二格。级数
	\begin{equation}\label{g4:eq:0405:wp-series}
		\wp_\Lambda(z)=\frac1{z^2}
		+\sum_{\lambda\in\Lambda\setminus\{0\}}
		\left(\frac1{(z-\lambda)^2}-\frac1{\lambda^2}\right)
	\end{equation}
	在 \(\mathbb C\setminus\Lambda\) 上正常收敛，即在每个紧集上各项绝对值的上确界之和有限。所得 Weierstrass 函数为偶亚纯函数，恰在各格点有主部 \(1/(z-\lambda)^2\)，并满足
	\[
		\wp_\Lambda(z+\lambda)=\wp_\Lambda(z)\quad(\lambda\in\Lambda),
		\qquad
		\wp_\Lambda'(z)=-2\sum_{\lambda\in\Lambda}(z-\lambda)^{-3}.
	\]
	令
	\[
		G_{2k}(\Lambda)=\sum_{\lambda\ne0}\lambda^{-2k}\qquad(k\ge2).
	\]
	若 \(0<|z|<\min_{\lambda\ne0}|\lambda|\)，则
	\begin{equation}\label{g4:eq:0405:wp-laurent}
		\wp_\Lambda(z)
		=\frac1{z^2}+\sum_{n\ge1}(2n+1)G_{2n+2}(\Lambda)z^{2n}.
	\end{equation}
\end{theorem}
\begin{proof}
	取一组正向格基 \(\lambda_1,\lambda_2\)，先估计格点的数目。实线性映射
	\((s,t)\mapsto s\lambda_1+t\lambda_2\) 可逆，所以存在 \(c>0\)，使
	\[
		|m\lambda_1+n\lambda_2|\ge c\sqrt{m^2+n^2}.
	\]
	因此半径 \(R\) 圆盘内的格点数为 \(O((R+1)^2)\)。把非零格点按 \(2^j\le|\lambda|<2^{j+1}\) 分组，便得
	\[
		\sum_{\lambda\ne0}|\lambda|^{-s}<\infty\qquad(s>2);
	\]
	有限个模长较小的格点单独计入即可。

	在紧集 \(K\subset\mathbb C\setminus\Lambda\) 上取 \(|z|\le R\)。当 \(|\lambda|>2R\) 时，
	\[
		\left|\frac1{(z-\lambda)^2}-\frac1{\lambda^2}\right|
		=\frac{|2\lambda z-z^2|}{|\lambda|^2|z-\lambda|^2}
		\le \frac{C_R}{|\lambda|^3}.
	\]
	有限个其余项在 \(K\) 上有界，故级数正常收敛。在任一格点附近，单独取出相应的二阶极点项，其余和全纯，因此得到所述亚纯延拓与主部。\cref{g4:prop:0104:holomorphic-limit}的 Cauchy 公式论证允许逐项求导；导数级数也由 \(\sum|\lambda|^{-3}\) 控制，得到导数公式。

	把 \(\lambda\) 换成 \(-\lambda\)，式\eqref{g4:eq:0405:wp-series}表明 \(\wp_\Lambda\) 为偶函数。导数级数绝对收敛，所以对任何 \(\mu\in\Lambda\) 可以重排求和，得到
	\(\wp_\Lambda'(z+\mu)=\wp_\Lambda'(z)\)。于是
	\[
		C_\mu(z)=\wp_\Lambda(z+\mu)-\wp_\Lambda(z)
	\]
	的极点主部抵消，且导数为零，因而为常数。偶性又给出
	\(C_\mu(-z-\mu)=-C_\mu(z)\)，所以该常数为零。这就证明原函数本身具有全部格周期。

	最后在 \(|z|<|\lambda|\) 时展开
	\[
		\frac1{(z-\lambda)^2}
		=\sum_{k\ge0}(k+1)\frac{z^k}{\lambda^{k+2}}.
	\]
	对任意固定的 \(r<\min_{\lambda\ne0}|\lambda|\)，减去常数项后的双重级数在 \(|z|\le r\) 上绝对一致收敛；可以先对 \(k\) 求和，再对 \(\lambda\) 求和。奇数次幂的系数因 \(\lambda,-\lambda\) 配对而消失，剩下的正偶数次项就是式\eqref{g4:eq:0405:wp-laurent}。
\end{proof}

为简化记号，在格固定时写 \(\wp=\wp_\Lambda\)。它在原点的前几项为
\[
	\wp(z)=z^{-2}+3G_4z^2+5G_6z^4+O(z^6).
\]
二阶极点与三阶极点分别出现在 \(\wp\) 和 \(\wp'\) 中，因此它们应满足三次关系。关系的系数可以只靠消去主部求出。

\begin{proposition}\label{g4:prop:0405:differential-equation}
	设 \(\Lambda\) 为秩二复格，记 \(\wp=\wp_\Lambda\)，并定义
	\begin{equation}\label{g4:eq:0405:invariants}
		g_2(\Lambda)=60G_4(\Lambda),\qquad
		g_3(\Lambda)=140G_6(\Lambda).
	\end{equation}
	则
	\begin{equation}\label{g4:eq:0405:wp-equation}
		(\wp')^2=4\wp^3-g_2\wp-g_3,
		\qquad
		\wp''=6\wp^2-\frac{g_2}{2}.
	\end{equation}
\end{proposition}
\begin{proof}
	写 \(a=3G_4\)、\(b=5G_6\)。由 Laurent 展开直接相乘，
	\[
		(\wp')^2-4\wp^3=-20a\,z^{-2}-28b+O(z^2).
	\]
	因此椭圆函数
	\[
		H=(\wp')^2-4\wp^3+g_2\wp+g_3
	\]
	在原点可去，且值为零。周期性使它在每个格点同样可去；其余地方本来全纯。所以 \(H\) 是整的椭圆函数，必为常数，进而恒为零。对第一式求导，在 \(\wp'\ne0\) 的开集上约去 \(2\wp'\)，得到第二式；两边都是亚纯函数，由恒等定理延拓到全部定义域。
\end{proof}

\subsection{复环面到光滑三次曲线}
\label{g4:subsec:0405:02}

微分方程已经给出一条三次曲线。还需核对它没有奇点，并判断 \(\wp,\wp'\) 是否遗漏或重复表示曲线上的点。关键是 \(\wp\) 在环面上只有一个二阶极点，因而到球面的次数恰为 \(2\)。

\begin{theorem}\label{g4:thm:0405:cubic-uniformization}
	设 \(\Lambda\subset\mathbb C\) 为秩二格，记
	\[
		\Delta(\Lambda)=g_2(\Lambda)^3-27g_3(\Lambda)^2.
	\]
	则 \(\Delta(\Lambda)\ne0\)，且
	\[
		E_\Lambda:\quad Y^2T=4X_1^3-g_2X_1T^2-g_3T^3
	\]
	为光滑射影三次曲线。映射
	\begin{equation}\label{g4:eq:0405:wp-uniformization}
		\Phi_\Lambda:T_\Lambda\longrightarrow E_\Lambda,\qquad
		[z]\longmapsto
		\begin{cases}
			[\wp(z):\wp'(z):1],&[z]\ne[0],\\
			[0:1:0],&[z]=[0],
		\end{cases}
	\end{equation}
	为双全纯同构。在仿射坐标 \(x=X_1/T,\ y=Y/T\) 中，
	\[
		\Phi_\Lambda^*\left(\frac{\mathrm{d}x}{y}\right)=\mathrm{d}z.
	\]
\end{theorem}
\begin{proof}
	\(\wp:T_\Lambda\to\mathbb P^1\) 唯一的极点是 \([0]\)，阶数为 \(2\)。由\cref{g4:thm:0202:degree}，此映射的次数为 \(2\)。偶性表明 \([z]\) 与 \([-z]\) 在同一纤维；若它们不同，这两个点已占满纤维。

	环面中恰有三个非零的二阶点
	\[
		h_1=[\lambda_1/2],\qquad
		h_2=[\lambda_2/2],\qquad
		h_3=[(\lambda_1+\lambda_2)/2].
	\]
	因为 \(\wp'\) 为奇函数且有格周期，\(\wp'(h_i)=0\)。它在 \([0]\) 有三阶极点，且没有其他极点。\cref{g4:prop:0203:principal-degree}使其零点总阶数等于 \(3\)，所以这三个零点恰好各为一阶，且没有其他零点。特别地，\(\wp''(h_i)\ne0\)，\(\wp\) 在这些点的局部映射次数均为 \(2\)。

	令 \(e_i=\wp(h_i)\)。若 \(e_i=e_j\)、\(i\ne j\)，同一纤维的重数至少为 \(4\)，与总次数为 \(2\) 矛盾。因此三个 \(e_i\) 互异。由式\eqref{g4:eq:0405:wp-equation}，
	\[
		4x^3-g_2x-g_3=4(x-e_1)(x-e_2)(x-e_3).
	\]
	三次多项式 \(4x^3-g_2x-g_3\) 有重根当且仅当 \(g_2^3=27g_3^2\)：重根 \(e\) 必满足 \(g_2=12e^2,\ g_3=-8e^3\)；反之，若这一系数等式成立，\(g_2=0\) 时也有 \(g_3=0\)，而 \(g_2\ne0\) 时取 \(e=-3g_3/(2g_2)\) 即得重根。故 \(\Delta\ne0\)。\cref{g4:prop:0204:cubic}于是保证 \(E_\Lambda\) 光滑、紧且连通。

	在 \([0]\) 附近，把齐次坐标同乘 \(z^3\)，得到
	\[
		[\wp(z):\wp'(z):1]
		=[z+O(z^5):-2+O(z^4):z^3].
	\]
	所以映射全纯延拓到所述无穷远点，并且 \(X_1/Y=-z/2+O(z^5)\) 是那里的一阶非退化坐标。其余点处，若 \(\wp'\ne0\)，第一坐标的导数非零；若 \(\wp'=0\)，则 \(\wp''\ne0\)，第二坐标的导数非零。因此 \(\Phi_\Lambda\) 处处为全纯浸入。

	若两点有相同像，先比较 \(\wp\) 坐标。次数为 \(2\) 和偶性迫使它们为 \([z]\)、\([-z]\)，或为同一点；在非二阶点，\(\wp'\) 在这两个点取相反的非零值，所以二者的像不同。无穷远点只有 \([0]\) 一个原像。故 \(\Phi_\Lambda\) 单射。它作为紧曲面间的非恒定全纯映射，像既开又闭，遂满射；各点的局部全纯逆拼成全局逆。

	在 \(\wp'\ne0\) 处，直接计算
	\(\mathrm{d}\wp/\wp'=\mathrm{d}z\)。左边在整个三次曲线上是全纯微分，有限分支点及无穷远处的可去性已在\cref{g4:prop:0204:cubic}中算出。因此此等式在全部环面上成立。
\end{proof}

反方向也能从积分得到。设 \(X\) 为亏格一的紧连通黎曼曲面，指定 \(O\in X\)，取非零全纯微分 \(\omega\)。它的零除子次数为 \(2g-2=0\)，所以处处非零。周期子群
\[
	\Lambda_\omega=\left\{\int_\gamma\omega:
		[\gamma]\in H_1(X,\mathbb Z)\right\}\subset\mathbb C
\]
由\cref{g4:prop:0401:period-lattice}为秩二格；\cref{g4:cor:0403:abel-embedding}给出的亏格一结论，在把归一化微分乘回原常数后，正是双全纯映射
\[
	A_O:X\longrightarrow\mathbb C/\Lambda_\omega,\qquad
	p\longmapsto\left[\int_O^p\omega\right].
\]
与 \(\Phi_{\Lambda_\omega}\) 复合，就得到 \(X\) 的三次模型。换一条非零全纯微分只会把格同乘一个非零复数，所选基点则确定环面的原点。

\ProseParagraphBreak

因此，带一个指定点的亏格一紧黎曼曲面、带原点的复环面和上述带无穷远点的光滑三次曲线相互对应。这样的带点曲线称为\term[tuoyuanquxian]{椭圆曲线}[elliptic curve]。指定点将在下一小节成为群的单位元。

\subsection{加法公式与几何群律}
\label{g4:subsec:0405:03}

复环面的加法由 \([u]+[v]=[u+v]\) 给出。要在三次曲线上计算它，需要把 \(\wp(u+v)\) 写成 \(\wp(u),\wp'(u),\wp(v),\wp'(v)\) 的有理式。消去极点的方法再次有效：固定 \(v\)，只要找出具有与 \(\wp(u+v)\) 相同主部的椭圆函数，两者就只能相差常数。

\begin{theorem}\label{g4:thm:0405:addition}
	设 \(\Lambda\) 为秩二复格，\(\wp=\wp_\Lambda\)。若 \(u,v\notin\Lambda\) 且 \(\wp(u)\ne\wp(v)\)，则
	\begin{equation}\label{g4:eq:0405:addition}
		\wp(u+v)
		=-\wp(u)-\wp(v)
		+\frac14
		\left(\frac{\wp'(u)-\wp'(v)}{\wp(u)-\wp(v)}\right)^2.
	\end{equation}
	若 \(u\notin\Lambda\) 且 \(\wp'(u)\ne0\)，则倍点公式为
	\begin{equation}\label{g4:eq:0405:duplication}
		\wp(2u)=-2\wp(u)
		+\frac14\left(\frac{\wp''(u)}{\wp'(u)}\right)^2.
	\end{equation}
	这些等式也可作为亚纯恒等式理解，在原表达出现可去奇点处取极限。
\end{theorem}
\begin{proof}
	先取 \(\wp'(v)\ne0\)，记 \(a=\wp(v)\)、\(b=\wp'(v)\)，并令 \(R(u)\) 为式\eqref{g4:eq:0405:addition}右边。由上一小节的二次纤维描述，\(\wp(u)-a\) 在环面上的零点恰为 \([v]\) 与 \([-v]\)，且均为简单零点。因此 \(R\) 只有在 \([0],[v],[-v]\) 才可能有极点。

	在 \(u=v\) 附近，分子与分母都消失，而分母的一阶系数为 \(b\ne0\)，故商可去。在 \(u=-v+t\) 附近，写 \(c=\wp''(v)\)，则
	\[
		\wp(u)-a=-bt+\frac c2t^2+O(t^3),\qquad
		\wp'(u)-b=-2b+ct+O(t^2).
	\]
	相除后常数项抵消，得到
	\[
		\frac{\wp'(u)-b}{\wp(u)-a}=\frac2t+O(t).
	\]
	因此 \(R\) 在 \([-v]\) 的主部恰为 \(t^{-2}\)，没有一阶极点项，与 \(\wp(u+v)\) 相同。

	在 \(u=0\) 附近，由 Laurent 展开，
	\[
		\frac{\wp'(u)-b}{\wp(u)-a}
		=-\frac2u-2au-bu^2+O(u^3).
	\]
	平方后再减去 \(\wp(u)+a\)，得到 \(R(u)=a+bu+O(u^2)\)。所以原点的极点可去，且 \(R(0)=a\)。于是 \(R(u)-\wp(u+v)\) 为整的椭圆函数，必为常数；在 \(u=0\) 取值，这个常数为零。加法公式对上述 \(v\) 成立。

	当 \(v\) 趋于一个非零二阶点时，只要 \(\wp(u)\ne\wp(v)\)，公式两边均全纯连续，因此等式延续到这一情形。这覆盖了陈述中的全部 \(u,v\)。最后令 \(v\to u\)，在 \(\wp'(u)\ne0\) 时差商趋于 \(\wp''(u)/\wp'(u)\)，即得倍点公式。
\end{proof}

令 \(P=(x_1,y_1)=\Phi_\Lambda([u])\)、\(Q=(x_2,y_2)=\Phi_\Lambda([v])\)。若 \(x_1\ne x_2\)，连接两点的直线斜率为
\[
	m=\frac{y_2-y_1}{x_2-x_1}.
\]
把 \(y=m(x-x_1)+y_1\) 代入 \(y^2=4x^3-g_2x-g_3\)，所得三次方程的三个根之和为 \(m^2/4\)。因此直线的第三个交点 \(R=(x_3,y_3)\)，按交点重数计，满足
\[
	x_3=\frac{m^2}{4}-x_1-x_2,\qquad
	y_3=m(x_3-x_1)+y_1.
\]
加法公式正好给出 \(\wp(u+v)=x_3\)。下面确定第二坐标的符号，并把所有重合情形一起纳入。

\begin{proposition}\label{g4:prop:0405:chord-tangent}
	设 \(\Lambda\) 为秩二复格，\(E_\Lambda\) 为\cref{g4:thm:0405:cubic-uniformization}中相应的光滑三次曲线，记 \(O=[0:1:0]\)。通过 \(\Phi_\Lambda\) 把 \(T_\Lambda\) 的加法搬到 \(E_\Lambda\)，得到以 \(O\) 为单位元的全纯交换群结构。逆元为
	\[
		-(x,y)=(x,-y),\qquad -O=O.
	\]
	若通过 \(P,Q\) 的射影直线与曲线的第三个交点为 \(R\)，重合时按切线和交点重数解释，则
	\[
		P+Q=-R.
	\]
	具体地，若 \(P=(x_1,y_1)\)、\(Q=(x_2,y_2)\) 均在仿射部分，且 \(x_1\ne x_2\)，记 \(m=(y_2-y_1)/(x_2-x_1)\)，则
	\[
		\begin{aligned}
			x(P+Q)&=\frac{m^2}{4}-x_1-x_2,\\
			y(P+Q)&=-y_1-m\bigl(x(P+Q)-x_1\bigr).
		\end{aligned}
	\]
	当 \(P=Q=(x,y)\)、\(y\ne0\) 时，用切线斜率
	\[
		m=\frac{12x^2-g_2}{2y}
	\]
	代入同一公式，并取 \(x_1=x_2=x\)。若 \(Q=-P\)，则 \(P+Q=O\)，包括 \(P=Q\) 为非零二阶点的情形。
\end{proposition}
\begin{proof}
	由双全纯同构搬运的运算满足全部群公理；其结合律直接来自复数加法，运算及取逆也全纯。偶函数 \(\wp\) 与奇函数 \(\wp'\) 给出所述逆元。

	先取非竖直直线 \(y=mx+b\)。函数 \(\ell=y-mx-b\) 在 \(O\) 的极点阶数恰为 \(3\)，因为 \(y\) 的极点阶数为 \(3\)，而 \(x\) 只有二阶极点。因此其零除子为三个交点之和，按零点阶数计，
	\[
		\operatorname{div}(\ell)=[P]+[Q]+[R]-3[O].
	\]
	这里局部零点阶数就是直线与光滑曲线的交点重数；切线使该阶数至少为 \(2\)。由\cref{g4:prop:0403:principal-periods}，主除子的 Abel--Jacobi 积分为零。在 \(T_\Lambda\) 的 \(\mathrm{d}z\) 周期坐标中，这句话恰为
	\[
		u+v+w\in\Lambda,
	\]
	其中 \(\Phi_\Lambda([w])=R\)。故 \(R=-(P+Q)\)，其第二坐标必须取反，得到斜线公式。若 \(P=Q\) 且 \(y\ne0\)，对三次方程隐式求导得
	\(2y\,\mathrm{d}y=(12x^2-g_2)\,\mathrm{d}x\)，因此切线斜率正如陈述。

	竖直线 \(x=a\) 的有限交点是一对 \((a,y),(a,-y)\)，\(y=0\) 时二者重合，另一个射影交点是 \(O\)。于是弦切规则给出互为逆元的两点之和为 \(O\)。通过 \(O\) 和任意有限点的直线也正是竖直线，所以同一规则给出 \(O+P=P\)。最后，\(O\) 处的切线为 \(T=0\)；局部参数表达中的 \(T\) 有三阶零点，因此三个交点都为 \(O\)，与 \(O+O=O\) 一致。所有情形均与已经定义的群法相符。
\end{proof}

弦切作图由此成为环面加法在三次模型中的表达，结合律也随同这一对应得到。第三章的零极点计数负责确定交点重数，Abel 定理则负责把交点之和转成群中的零。

\subsection{格变换、j 不变量与退化}
\label{g4:subsec:0405:04}

先区分两种改变：重新选择同一个格的整数基不会改变格，把格同乘一个非零复数则改变其在复平面中的大小和方向，但不改变复环面的同构类。实际上，这已经给出了全部同构。

\begin{proposition}\label{g4:prop:0405:lattice-isomorphism}
	设 \(\Lambda,\Lambda'\subset\mathbb C\) 为秩二格。保持原点的双全纯映射
	\[
		F:\mathbb C/\Lambda\longrightarrow\mathbb C/\Lambda'
	\]
	恰为 \(F([z])=[az]\)，其中 \(a\in\mathbb C^\times\) 满足 \(a\Lambda=\Lambda'\)。任意不要求保持原点的双全纯映射都是这种映射再复合一次目标平移。
\end{proposition}
\begin{proof}
	目标环面有处处非零的微分 \(\mathrm{d}w\)。其拉回与源上的 \(\mathrm{d}z\) 相比，商是紧连通曲面上的全纯函数，所以
	\[
		F^*(\mathrm{d}w)=a\,\mathrm{d}z
	\]
	其中 \(a\) 为常数。双全纯性保证 \(a\ne0\)。沿源上由格向量 \(\lambda\) 给出的闭路积分，得到 \(a\lambda\in\Lambda'\)，因而 \(a\Lambda\subset\Lambda'\)。积分从 \(0\) 到 \(z\) 的道路，并使用 \(F([0])=[0]\)，便得 \(F([z])=[az]\)：局部积分值的差是目标周期，因此在商中这一等式不依赖道路。对 \(F^{-1}\) 作同样论证，得到 \(\Lambda'\subset a\Lambda\)，故二者相等。

	反过来，\(a\Lambda=\Lambda'\) 时，线性映射 \(z\mapsto az\) 及其逆都下降到商空间，给出所述双全纯映射。若原映射不保持原点，先在目标减去 \(F([0])\) 即可应用已经证明的情形。
\end{proof}

把格写成 \(\Lambda=\lambda_1(\mathbb Z+\tau\mathbb Z)\)，其中 \(\operatorname{Im}\tau>0\)，便去掉了整体缩放。令 \(\mathbb H=\{\tau\in\mathbb C:\operatorname{Im}\tau>0\}\)。同一个格的另一组正向整数基可以写成
\[
	\lambda_1'=d\lambda_1+c\lambda_2,\qquad
	\lambda_2'=b\lambda_1+a\lambda_2,\qquad
	\begin{pmatrix}a&b\\c&d\end{pmatrix}\in\operatorname{SL}(2,\mathbb Z).
\]
于是
\begin{equation}\label{g4:eq:0405:modular-action}
	\tau'=\frac{a\tau+b}{c\tau+d},\qquad
	\operatorname{Im}\tau'=\frac{\operatorname{Im}\tau}{|c\tau+d|^2}>0.
\end{equation}
反过来，同构的复环面对应相差一个非零复数倍的格，把两边的正向基比较，就得到这样的整数矩阵。因此复环面的同构类恰由 \(\mathbb H\) 中这些变换的轨道参数化。这里正向性使整数行列式为 \(1\)，因为复数乘法保持实平面的定向。

\ProseParagraphBreak

格级数还直接给出缩放公式
\begin{equation}\label{g4:eq:0405:homogeneity}
	\wp_{a\Lambda}(az)=a^{-2}\wp_\Lambda(z),\qquad
	g_2(a\Lambda)=a^{-4}g_2(\Lambda),\qquad
	g_3(a\Lambda)=a^{-6}g_3(\Lambda).
\end{equation}
在式\eqref{g4:eq:0405:wp-series}中令格点为 \(a\lambda\) 即可逐项验证。因此
\begin{equation}\label{g4:eq:0405:j}
	j(\Lambda)\coloneqq
	1728\,\frac{g_2(\Lambda)^3}{g_2(\Lambda)^3-27g_3(\Lambda)^2}
\end{equation}
不受缩放影响，定义了复环面以及相应椭圆曲线的\term[jbubianliang]{\(j\) 不变量}[\(j\)-invariant]。常数 \(1728\) 是此处采用的归一化。

\begin{theorem}\label{g4:thm:0405:j-classification}
	两个一维复环面双全纯同构，当且仅当它们的 \(j\) 不变量相同；若各自指定原点，可以取保持原点的同构。等价地，对无重根的三次模型
	\[
		E_{A,B}:y^2=4x^3-Ax-B,\qquad
		j(E_{A,B})=1728\frac{A^3}{A^3-27B^2},
	\]
	两条曲线有保持无穷远点的双全纯同构，当且仅当相应 \(j\) 值相同。因而这个数也与同一带点亏格一曲面所选的三次模型无关。
\end{theorem}
\begin{proof}
	先核对三次模型之间的全部带点同构。设
	\(F:(E_{A,B},O)\to(E_{A',B'},O')\) 双全纯。两条曲线均有亏格 \(1\)，其坐标 \(x,y\) 在 \(O\) 的极点阶数分别为 \(2,3\)。由式\eqref{g4:eq:0305:elliptic-pole-count}，
	\[
		H^0(\mathcal O(2[O]))=\operatorname{span}\{1,x\},\qquad
		H^0(\mathcal O(3[O]))=\operatorname{span}\{1,x,y\}.
	\]
	拉回保持极点的精确阶数，所以
	\[
		F^*x'=\alpha x+\beta,\qquad
		F^*y'=\gamma y+\delta x+\varepsilon,\qquad
		\alpha\gamma\ne0.
	\]
	代入目标方程，再用 \(y^2=4x^3-Ax-B\)。六个函数
	\(1,x,y,x^2,xy,x^3\) 的极点阶数互异，线性无关；比较 \(xy,y,x^2\) 的系数，依次得到
	\(\delta=\varepsilon=\beta=0\)。比较 \(x^3,x,1\) 的系数，得到
	\[
		\gamma^2=\alpha^3,\qquad
		A'=\alpha^2 A,\qquad B'=\alpha^3 B.
	\]
	令 \(u=\gamma/\alpha\)，便有 \(u^2=\alpha,\ u^3=\gamma\)。所以所有这样的同构都按
	\[
		(x,y)\longmapsto(u^2x,u^3y),\qquad
		(A',B')=(u^4A,u^6B)
	\]
	作用，显然保持 \(j\)。

	反过来，假设 \(j(E_{A,B})=j(E_{A',B'})\)。由于两边判别式都非零，交叉相乘得到
	\[
		A^3(B')^2=(A')^3B^2.
	\]
	若 \(A=0\)，则 \(A'=0\)，且 \(B,B'\ne0\)，取 \(u^6=B'/B\) 即可。若 \(B=0\)，同理 \(B'=0\)，取 \(u^4=A'/A\) 即可。剩余情形中四个系数都非零，记
	\[
		r=A'/A,\qquad s=B'/B.
	\]
	有 \(r^3=s^2\)。取 \(u\) 满足 \(u^2=s/r\)，则 \(u^4=r,\ u^6=s\)，因此同样得到所需坐标同构。这证明了三次模型的等价判据。

	\(\Phi_\Lambda\) 将各复环面识别为 \(E_{g_2(\Lambda),g_3(\Lambda)}\)，故相同 \(j\) 给出保持原点的环面同构。另一方面，任意环面同构先减去目标中的像点后保持原点，\cref{g4:prop:0405:lattice-isomorphism}使两格相差缩放，式\eqref{g4:eq:0405:homogeneity}于是保证其 \(j\) 相同。任意带点亏格一曲面又由 Abel--Jacobi 积分同构于这样的环面，所以两种判据完全相容。
\end{proof}

\begin{example}\label{g4:exm:0405:symmetric-lattices}
	方格 \(\Lambda=\mathbb Z+\mathrm{i}\mathbb Z\) 满足 \(\mathrm{i}\Lambda=\Lambda\)。缩放公式给出
	\(g_3=\mathrm{i}^{-6}g_3=-g_3\)，故 \(g_3=0\)；判别式非零使 \(g_2\ne0\)，因此 \(j=1728\)。若
	\(\zeta=\exp(\pi\mathrm{i}/3)\)、\(\Lambda=\mathbb Z+\zeta\mathbb Z\)，则 \(\zeta^2=\zeta-1\)，故 \(\zeta\Lambda=\Lambda\)。这时 \(g_2=\zeta^{-4}g_2\) 迫使 \(g_2=0\)，而 \(g_3\ne0\)，因此 \(j=0\)。格的额外旋转对称性，分别表现为三次方程中一个系数消失。
\end{example}

每个复数都能作为这里的 \(j\) 值出现。给定 \(j_0\ne0\)，取
\[
	A=3,\qquad B^2=1-\frac{1728}{j_0};
\]
则 \(A^3-27B^2=27\cdot1728/j_0\ne0\)，相应三次曲线的 \(j\) 恰为 \(j_0\)。\(j_0=0\) 时取 \(A=0,B=1\)。这些光滑三次曲线又都由前面的积分构造来自格。因此 \(j\) 在同构类的集合上给出完整的复数参数。

\ProseParagraphBreak

判别式为零时，三次模型失去光滑性。若 \((A,B)\ne(0,0)\) 且 \(A^3=27B^2\)，存在 \(e\ne0\)，使
\[
	4x^3-Ax-B=4(x-e)^2(x+2e).
\]
在奇点 \((e,0)\) 附近令 \(u=x-e\)，方程变成
\[
	y^2=4u^2(u+3e).
\]
取全纯函数 \(h(u)=2\sqrt{u+3e}\)，其中 \(h(0)\ne0\)，再令
\[
	U=y-u h(u),\qquad V=y+u h(u).
\]
这确为局部可逆的全纯坐标变换：由 \(U,V\) 先恢复 \(y=(U+V)/2\)，再对导数为 \(h(0)\ne0\) 的一元函数 \(u\mapsto u h(u)\) 求局部逆。原方程遂成为 \(UV=0\)，即两条光滑复分支横截相遇的\term[jiedian]{结点}[node]。若 \(A=B=0\)，方程为
\[
	y^2=4x^3,\qquad x=t^2,\quad y=2t^3;
\]
这是只有一条分支的\term[jiandian]{尖点}[cusp]模型，其参数化在 \(t=0\) 的导数为零。两种退化的区别可以直接由局部方程读出。

\ProseParagraphBreak

格的一个周期趋于无穷远时，会出现结点退化。保留周期 \(1\)，令另一个周期 \(\tau\) 的虚部趋于无穷大，则远离实轴的格行逐渐退出任意固定有界区域；相应级数剩下整数格上的一维和。

\begin{proposition}\label{g4:prop:0405:trigonometric-limit}
	设 \(\Lambda_\tau=\mathbb Z+\tau\mathbb Z\)，其中 \(\operatorname{Im}\tau>0\)。当 \(\operatorname{Im}\tau\to+\infty\) 时，在 \(\mathbb C\setminus\mathbb Z\) 的每个紧集上有一致极限
	\begin{equation}\label{g4:eq:0405:trigonometric-limit}
		\wp_{\Lambda_\tau}(z)\longrightarrow
		W(z)\coloneqq\pi^2\csc^2(\pi z)-\frac{\pi^2}{3}.
	\end{equation}
	相应不变量满足
	\[
		g_2(\Lambda_\tau)\longrightarrow\frac{4\pi^4}{3},
		\qquad
		g_3(\Lambda_\tau)\longrightarrow\frac{8\pi^6}{27},
		\qquad |j(\Lambda_\tau)|\longrightarrow\infty.
	\]
	极限三次方程为
	\[
		y^2=4\left(x+\frac{\pi^2}{3}\right)^2
			\left(x-\frac{2\pi^2}{3}\right),
	\]
	在 \((-\pi^2/3,0)\) 有一个结点。
\end{proposition}
\begin{proof}
	记 \(t=\operatorname{Im}\tau\)，固定紧集 \(K\subset\mathbb C\setminus\mathbb Z\)，并取 \(|z|\le R\)。当 \(t>2R+1\) 时，所有 \(n\ne0\) 的格点 \(m+n\tau\) 都在 \(K\) 之外，且可用之前的估计：
	\[
		\left|\frac1{(z-m-n\tau)^2}-\frac1{(m+n\tau)^2}\right|
		\le C_R|m+n\tau|^{-3}.
	\]
	对任意实数 \(a\) 及 \(h\ge1\)，有
	\[
		\sum_{m\in\mathbb Z}
		\bigl((m+a)^2+h^2\bigr)^{-3/2}
		\le C(h^{-3}+h^{-2})\le C' h^{-2}.
	\]
	这个界与 \(a\) 无关：先取离 \(-a\) 最近的整数，再把其余整数向两侧编号，并把递减函数的和与积分
	\[
		\int_0^\infty(s^2+h^2)^{-3/2}\,\mathrm{d}s=h^{-2}
	\]
	比较即可。将 \(a=n\operatorname{Re}\tau\)、\(h=|n|t\) 代入，再对 \(n\ne0\) 求和，得到
	\[
		\sum_{n\ne0}\sum_{m\in\mathbb Z}|m+n\tau|^{-3}
		\le\frac{C''}{t^2}\sum_{n\ne0}\frac1{n^2}
		=O(t^{-2}).
	\]
	所以非零格行对 \(\wp\) 的总贡献在 \(K\) 上一致趋于零，剩下
	\[
		W(z)=\frac1{z^2}
		+\sum_{m\in\mathbb Z\setminus\{0\}}
		\left((z-m)^{-2}-m^{-2}\right).
	\]

	现计算这个一维和。级数
	\(S(z)=\sum_{m\in\mathbb Z}(z-m)^{-2}\) 正常收敛并以 \(1\) 为周期。它与 \(\pi^2\csc^2(\pi z)\) 在每个整数处有相同主部，因此两者之差整且以 \(1\) 为周期。在竖条 \(0\le\operatorname{Re}z\le1\) 上，当 \(|\operatorname{Im}z|\to\infty\) 时，积分比较给出
	\[
		|S(z)|\le \sum_m|z-m|^{-2}
		=O(|\operatorname{Im}z|^{-1})\longrightarrow0,
	\]
	而 \(\csc^2(\pi z)\) 指数趋于零。差函数在竖条中间的紧部分也有界；周期性使差函数在全平面有界，前述 Cauchy 估计遂使它为常数，再由上述极限知常数为零。于是 \(S(z)=\pi^2\csc^2(\pi z)\)。在原点比较常数项，
	\[
		\sum_{m\ne0}m^{-2}=\frac{\pi^2}{3},
	\]
	即得式\eqref{g4:eq:0405:trigonometric-limit}。

	取固定圆周 \(|z|=r<1/2\)。已经证明的收敛在该圆周上一致，而去掉共同的 \(z^{-2}\) 后，各函数在圆内全纯；Cauchy 系数公式因而允许对 Laurent 系数取极限。由正弦的 Taylor 展开，
	\[
		W(z)=z^{-2}+\frac{\pi^4}{15}z^2
			+\frac{2\pi^6}{189}z^4+O(z^6).
	\]
	式\eqref{g4:eq:0405:invariants}与\eqref{g4:eq:0405:wp-laurent}分别把这两个系数乘以 \(20\) 和 \(28\)，得到所述 \(g_2,g_3\) 极限。于是
	\[
		g_2^3-27g_3^2\longrightarrow0,\qquad
		g_2^3\longrightarrow\frac{64\pi^{12}}{27}\ne0.
	\]
	每个有限 \(t\) 对应真正的秩二格，判别式仍非零，所以 \(j\) 的模趋于无穷大。代入极限系数并因式分解，便得到所述结点方程。
\end{proof}

还可看清结点处的两个分支来自哪里。在极限圆柱 \(\mathbb C/\mathbb Z\) 上取 \(q=\exp(2\pi\mathrm{i}z)\)，它把圆柱双全纯地识别为 \(\mathbb C^\times\)，且
\[
	W(z)=-\frac{4\pi^2q}{(q-1)^2}-\frac{\pi^2}{3}.
\]
圆柱的两端 \(q\to0\)、\(q\to\infty\) 都趋于同一个值 \(-\pi^2/3\)，并且 \(W'(z)\to0\)。令 \(u=W+\pi^2/3\)、\(y=W'\)，在 \(q\to0\) 与 \(q\to\infty\) 时分别有 \(y/u\to2\pi\mathrm{i}\) 与 \(y/u\to-2\pi\mathrm{i}\)，所以两端确实落在同一结点的两条不同分支上。这样，一个周期无限增长的解析极限，与三次方程中两个根合并的代数极限给出了相同的几何变化。
% 纲要条目排为编号小节标题，并进入目录与书签；说明段仍保留于纲要注释中。


\begin{chaptersummary}
	周期把局部全纯微分与曲面的闭路联系起来。双线性公式使 \(a\) 周期可唯一归一化，并给出 \(\tau^{\mathsf T}=\tau\)、\(\operatorname{Im}\tau>0\)，所以 \(\mathbb Z^g+\tau\mathbb Z^g\) 确为格。点的积分模去这个格后成为单值映射；对称积上的求值与喷射计算则把映射的微分秩转化为全纯微分的消失条件。

	\ProseParagraphBreak

	对度零除子构造第三类微分后，带极点互反公式把其剩余周期写成除子的 Abel--Jacobi 积分。积分类为零恰好允许减去一条全纯微分，使全部周期为 \(2\pi\mathrm{i}\) 的整数倍；取指数便恢复具有指定除子的亚纯函数。反演进一步说明所有积分值都能实现，同一积分值的有效除子组成完全线性系，其维数由 Riemann--Roch 给出。

	\ProseParagraphBreak

	theta 级数的变换因子定义全纯线丛，其正曲率代表标准主极化。本章直接证明了这些解析与线性代数结论；将 theta 零除子识别为 \(W_{g-1}\) 的平移，以及从主极化 Jacobi 簇重构曲面，分别使用文献中的 Riemann theta 定理与 Torelli 定理。亏格一的情形则由格级数完整算出：\(\wp,\wp'\) 给出三次模型，环面加法给出弦切群律，格的复数缩放给出 \(j\) 分类，而三角极限展示光滑三次曲线如何趋于结点曲线。
\end{chaptersummary}

\BookExercises

\begin{exercise}
	\label{g4:ex:0401:exchange-cycles}
	设 \(X\) 为亏格 \(g\ge1\) 的紧连通黎曼曲面，归一化周期矩阵为 \(\tau\)。交换同调辛基为 \(a'_i=b_i,\ b'_i=-a_i\)。求新的归一化微分与周期矩阵，直接验证新矩阵的虚部正定，并写出两种周期表示之间的复环面同构。
\end{exercise}
\begin{solution}
	\(\tau\) 可逆。事实上，若 \(\tau v=0\)，则
	\[
		0=\operatorname{Im}(\overline v^{\mathsf T}\tau v)
		=\overline v^{\mathsf T}(\operatorname{Im}\tau)v,
	\]
	正定性迫使 \(v=0\)。记旧微分列向量为 \(\boldsymbol\omega\)，取
	\[
		\omega'_j=\sum_k(\tau^{-1})_{kj}\omega_k,\qquad
		\boldsymbol\omega'=\tau^{-1}\boldsymbol\omega.
	\]
	后一式用了 \(\tau^{-1}\) 对称。其 \(a'\) 周期为 \(\tau\tau^{-1}=I\)，其 \(b'\) 周期为
	\[
		\tau'=-\tau^{-1}.
	\]
	这个矩阵仍对称。设 \(Y=\operatorname{Im}\tau\)，直接相减得到
	\[
		\operatorname{Im}(-\tau^{-1})
		=\overline\tau^{-1}Y\tau^{-1}.
	\]
	对任意非零列向量 \(v\)，右边的二次型等于 \(\overline{(\tau^{-1}v)}^{\mathsf T}Y(\tau^{-1}v)>0\)，故正定。最后
	\[
		\tau^{-1}(\mathbb Z^g+\tau\mathbb Z^g)
		=\mathbb Z^g+\tau^{-1}\mathbb Z^g
		=\mathbb Z^g+\tau'\mathbb Z^g.
	\]
	因此 \([z]\mapsto[\tau^{-1}z]\) 给出所求双全纯群同构；它正是同一积分泛函改用新归一化微分基后的坐标变换。
\end{solution}

\begin{exercise}
	\label{g4:ex:0403:second-kind}
	设 \(X\) 为亏格 \(g\ge1\) 的紧连通黎曼曲面，在 \(p\) 取中心坐标 \(z\)，给定 \(c\in\mathbb C\)。固定避开 \(p\) 的几何辛基及相应归一化微分 \(\omega_j\)。证明存在唯一亚纯微分 \(\eta\)，只在 \(p\) 可能有极点，主部为 \(c\,\mathrm{d}z/z^2\)，并且全部 \(a\) 周期为零。证明
	\[
		\int_{b_j}\eta
		=2\pi\mathrm{i}c\left(\frac{\omega_j}{\mathrm{d}z}\right)(p).
	\]
	这种留数全为零的亚纯微分称为第二类微分；这里主部连同所选坐标一起给定。
\end{exercise}
\begin{solution}
	对 \(L=K_X\) 使用\cref{g4:cor:0303:principal-parts-criterion}，对偶测试空间为常数，所需条件只是主部的留数为零。这对 \(c\,\mathrm{d}z/z^2\) 成立，因此存在具有该主部且没有其他极点的微分 \(\eta_0\)。减去 \(\sum_i(\int_{a_i}\eta_0)\omega_i\)，便使全部 \(a\) 周期为零。同主部的两个解之差全纯，若其 \(a\) 周期全为零，归一化微分的唯一性使这个差为零。

	在沿几何辛基切开的多边形内，取 \(\omega_j\) 的单值原函数 \(F_j\)。\cref{g4:prop:0403:reciprocity}的配对边计算不依赖极点阶数；删去 \(p\) 周围小圆盘后，边界公式仍给出
	\[
		\sum_i\left(\int_{a_i}\omega_j\int_{b_i}\eta
		-\int_{b_i}\omega_j\int_{a_i}\eta\right)
		=2\pi\mathrm{i}\operatorname{Res}_p(F_j\eta).
	\]
	区别仅在最后的留数。写
	\[
		F_j(z)=F_j(0)+a_jz+O(z^2),\qquad
		a_j=\left(\frac{\omega_j}{\mathrm{d}z}\right)(p),
	\]
	便有 \(\operatorname{Res}_p(F_j\eta)=ca_j\)。左边因 \(a\) 周期归一化而等于 \(\int_{b_j}\eta\)，得到所求公式。这里没有从主部直接读出微分的全部周期，而是由原函数的一阶项补出了相应信息。
\end{solution}

\begin{exercise}
	\label{g4:ex:0404:genus-three-fibers}
	设 \(X\) 为亏格 \(3\) 的紧连通黎曼曲面。若给定二重全纯映射 \(\pi:X\to\mathbb P^1\)，取其任一纤维除子 \(H\)，证明 \(A_q^{(2)}\) 在 \(A_q^{(2)}(H)\) 上的纤维恰由 \(\pi\) 的全部纤维除子组成，是一个嵌入的 \(\mathbb P^1\)，且 Abel 映射沿这个纤维的微分秩处处为 \(1\)。另证明：若 \(X\) 不是超椭圆曲面，则 \(A_q^{(2)}\) 的每个非空纤维都是一点，微分秩处处为 \(2\)。
\end{exercise}
\begin{solution}
	在靶球面选坐标 \(t\)，使 \(H=\pi^*[\infty]\)，并写 \(f=t\circ\pi\)。则 \(f\) 的极点除子恰为 \(H\)，所以 \(1,f\) 是 \(H^0(\mathcal O(H))\) 中的两个线性无关元素。由\cref{g4:thm:0305:canonical-map}对任意次数 \(2\) 的有效除子给出的界，\(\ell(H)\le2\)，因此
	\[
		H^0(\mathcal O(H))=\operatorname{span}\{1,f\}.
	\]
	对 \([a:b]\in\mathbb P^1\)，截面 \(a+bf\) 的有效零除子为 \(H+\operatorname{div}(a+bf)\)。当 \(b\ne0\) 时它是 \(\pi^*[-a/b]\)，当 \(b=0\) 时它是 \(H\)。\cref{g4:thm:0404:abel-fibers}因而说明整个 Abel 纤维正是这族纤维除子，并给出其全纯嵌入。此论证包括分支点处的重合除子 \(2[p]\)。

	对该族中任一 \(D\)，都有 \(\ell(D)=2\)。Riemann--Roch 在 \(g=3,d=2\) 时给出
	\[
		h^0(K_X(-D))=\ell(D)=2.
	\]
	故\cref{g4:prop:0402:abel-differential}给出秩 \(3-2=1\)。这也说明重合点必须在对称积坐标中计算，不能把两个相同点的求值列简单重复后就视为全部微分。

	若 \(X\) 不是超椭圆曲面而某个次数 \(2\) 的有效除子 \(D\) 有 \(\ell(D)\ge2\)，则存在非常值函数，其极点除子不超过 \(D\)。它到球面的次数为 \(1\) 或 \(2\)；次数 \(1\) 迫使 \(X\) 为球面，次数 \(2\) 与非超椭圆假设矛盾。因此所有这样的 \(D\) 都有 \(\ell(D)=1\)。纤维为 \(\mathbb P^0\)，而上述维数等式给出 \(h^0(K_X(-D))=1\)，故微分秩为 \(2\)。
\end{solution}

\begin{exercise}
	\label{g4:ex:0403:torsion-divisor}
	设 \(E=\mathbb C/\Lambda\)，\(O=[0]\)，\(n\ge1\)。证明 \(n[p]-n[O]\) 是主除子当且仅当 \(np=O\)，并证明这样的点恰有 \(n^2\) 个。若 \(p\ne O\) 满足此条件，求具有该除子的亚纯函数到球面的次数，并描述这些函数之间的差别。
\end{exercise}
\begin{solution}
	在 \(\mathrm{d}z\) 的积分坐标中，Abel--Jacobi 映射就是环面的恒等映射。因此
	\[
		\operatorname{AJ}(n[p]-n[O])=np,
	\]
	\cref{g4:thm:0403:abel}给出第一个等价。设 \(\Lambda=\mathbb Z\lambda_1+\mathbb Z\lambda_2\)，条件 \(nz\in\Lambda\) 等价于
	\[
		p=\left[\frac{a\lambda_1+b\lambda_2}{n}\right],
		\qquad 0\le a,b<n.
	\]
	这些 \(n^2\) 个类两两不同：两代表之差属于 \(\Lambda\) 时，格基的实线性无关性迫使两个整数差都被 \(n\) 整除，而它们的绝对值小于 \(n\)。当 \(p\ne O\) 时，所得函数恰在 \(O\) 有一个 \(n\) 阶极点，所以映射次数为 \(n\)。同除子的两函数之比在紧曲面上处处全纯且非零，故它们相差非零常数倍。
\end{solution}

\begin{exercise}
	\label{g4:ex:0405:flexes}
	设 \(E:y^2=4x^3-g_2x-g_3\) 为光滑三次曲线，单位元为无穷远点 \(O\)。把切线与曲线在该点的交点重数为 \(3\) 的点称为拐点。证明拐点恰为 \(3\) 扭点，因而共有九个。证明有限拐点都满足 \(y\ne0\)，其横坐标恰为
	\[
		48x^4-24g_2x^2-48g_3x-g_2^2=0
	\]
	的根，并说明这个四次多项式有四个不同的根。
\end{exercise}
\begin{solution}
	在 \(P\) 处的切线有交点除子 \(2[P]+[R]\)，其中允许 \(R=P\)。由弦切群律，\(2P+R=O\)。因此切线具有三重交点当且仅当 \(R=P\)，亦即 \(3P=O\)。无穷远处的切线 \(T=0\) 有三重交点，故 \(O\) 包括在内。由\cref{g4:ex:0403:torsion-divisor}，三扭点恰有 \(9\) 个。

	若有限点满足 \(y=0\)，它是非零二扭点，不能同时满足 \(3P=O\)，因为后者与 \(2P=O\) 会推出 \(P=O\)。所以有限拐点的 \(y\ne0\)。它的切线斜率为 \(m=(12x^2-g_2)/(2y)\)，第三个交点的横坐标为 \(m^2/4-2x\)。它等于 \(x\) 当且仅当
	\[
		(12x^2-g_2)^2=48xy^2=48x(4x^3-g_2x-g_3).
	\]
	移项便得到题中的四次式。由于切线不是竖直线，在同一横坐标处只能给出同一个纵坐标，所以这个横坐标条件也充分保证第三交点为 \(P\)。

	反过来，四次式的根不能同时使 \(4x^3-g_2x-g_3=0\)。否则上述等式会给出 \(12x^2-g_2=0\)，使三次多项式有重根，与光滑性矛盾。因此每个四次式的根给出两点 \((x,\pm y)\)，二者都是拐点。八个非零三扭点按互为逆元配成四对；每对有相同横坐标，二次映射 \(x:E\to\mathbb P^1\) 不容许两对共享一个横坐标。因此恰有四个不同的根。
\end{solution}
\begin{exercise}
	\label{g4:ex:0404:theta-powers}
	设 \(\tau=\tau^{\mathsf T}\in M_g(\mathbb C)\)、\(\operatorname{Im}\tau>0\)，\(A=\mathbb C^g/(\mathbb Z^g+\tau\mathbb Z^g)\)，\(L_\theta\) 为正文由 theta 乘子构造的线丛。对整数 \(N\ge1\)，证明
	\[
		h^0(A,L_\theta^{\otimes N})=N^g.
	\]
	更具体地，对 \(r\in\{0,\ldots,N-1\}^g\) 定义
	\[
		\Theta_r^{(N)}(z)
		=\sum_{k\in\mathbb Z^g}
		\exp\left(\frac{\pi\mathrm{i}}N(r+Nk)^{\mathsf T}
		\tau(r+Nk)+2\pi\mathrm{i}(r+Nk)^{\mathsf T}z\right).
	\]
	证明这些函数给出截面空间的一组基。
\end{exercise}
\begin{solution}
	\(L_\theta^{\otimes N}\) 的截面提升为整函数 \(f\)，满足
	\[
		f(z+m+\tau n)
		=e^{-N\pi\mathrm{i}n^{\mathsf T}\tau n
		-2N\pi\mathrm{i}n^{\mathsf T}z}f(z).
	\]
	整数周期性给出 Fourier 展开
	\[
		f(z)=\sum_{\ell\in\mathbb Z^g}c_\ell
		e^{2\pi\mathrm{i}\ell^{\mathsf T}z}.
	\]
	如\cref{g4:prop:0404:theta-section-dimension}的证明，这个展开来自逐变量 Laurent 展开，在紧集上局部一致收敛，且系数唯一。比较 \(\tau n\) 平移前后的系数，得到
	\[
		c_{\ell+Nn}
		=\exp\bigl(N\pi\mathrm{i}n^{\mathsf T}\tau n
		+2\pi\mathrm{i}\ell^{\mathsf T}\tau n\bigr)c_\ell.
	\]
	因此每个余类 \(\ell\bmod N\mathbb Z^g\) 中的全部系数由其中一个决定，截面空间维数至多为 \(N^g\)。

	\(\Theta_r^{(N)}\) 的指数实部在紧集上被某个 \(-c|k|^2+C|k|+C'\) 控制，其中 \(c>0\)。所以级数及其各阶导数正常收敛，定义整函数。整数平移使每项不变；对 \(\tau n\) 平移，把 \(r+Nk\) 改写为 \(r+N(k+n)-Nn\)，完成平方后恰得到上面的乘子。因此它们都是所需截面。不同 \(r\) 的 Fourier 支集位于不同余类，而每个余类上至少一个系数非零，所以它们线性无关。已有 \(N^g\) 个独立截面，与前述上界合起来便证明它们构成基。
\end{solution}

\begin{exercise}
	\label{g4:ex:0404:theta-characteristics}
	设 \(\tau=\tau^{\mathsf T}\)、\(\operatorname{Im}\tau>0\)，取 \(\alpha,\beta\in\{0,\tfrac12\}^g\)，定义带半整特征的 theta 函数
	\[
		\theta\left[\begin{matrix}\alpha\\ \beta\end{matrix}\right](z,\tau)
		=\sum_{n\in\mathbb Z^g}
		\exp\left(\pi\mathrm{i}(n+\alpha)^{\mathsf T}\tau(n+\alpha)
		+2\pi\mathrm{i}(n+\alpha)^{\mathsf T}(z+\beta)\right).
	\]
	证明它关于 \(z\) 为偶函数或奇函数，奇偶性由 \((2\alpha)^{\mathsf T}(2\beta)\bmod2\) 决定。计算奇特征的数目，并证明每个奇特征对应的函数在 \(z=0\) 为零。
\end{exercise}
\begin{solution}
	虚部正定仍给出高斯衰减，所以级数正常收敛，可以重排。由于 \(2\alpha\in\mathbb Z^g\)，在 \(-z\) 的级数中作指标替换 \(n\mapsto-n-2\alpha\)，得到
	\[
		\theta\left[\begin{matrix}\alpha\\ \beta\end{matrix}\right](-z,\tau)
		=e^{-4\pi\mathrm{i}\alpha^{\mathsf T}\beta}
		\theta\left[\begin{matrix}\alpha\\ \beta\end{matrix}\right](z,\tau).
	\]
	这里有关 \(n\) 的额外因子是 \(e^{-4\pi\mathrm{i}n^{\mathsf T}\beta}=1\)。常数因子等于 \((-1)^{(2\alpha)^{\mathsf T}(2\beta)}\)，因此得所述奇偶性。每个函数都不恒为零：乘以 \(e^{-2\pi\mathrm{i}\alpha^{\mathsf T}z}\) 后成为整数周期的 Fourier 级数，而其各系数都是非零指数。

	写 \(a=2\alpha,b=2\beta\in\{0,1\}^g\)，在二元域上计数。若 \(a=0\)，则没有 \(a^{\mathsf T}b=1\) 的 \(b\)。若 \(a\ne0\)，选 \(a_j=1\) 的一个坐标；其余 \(g-1\) 个 \(b\) 分量可任意选，而 \(b_j\) 被条件唯一确定。因此每个非零 \(a\) 对应 \(2^{g-1}\) 个奇特征，奇特征总数为
	\[
		(2^g-1)2^{g-1}.
	\]
	对奇函数在 \(z=0\) 代入 \(f(0)=-f(0)\)，便得 \(f(0)=0\)。这一结论没有断言偶特征在原点一定非零。
\end{solution}

\begin{exercise}
	\label{g4:ex:0404:residual-symmetry}
	设 \(X\) 是亏格 \(g\ge1\) 的紧连通黎曼曲面，固定 \(q\in X\)，并令
	\[
		W_{g-1}=A_q^{(g-1)}(\operatorname{Sym}^{g-1}X),\qquad
		k=\operatorname{AJ}(K-(2g-2)[q]),
	\]
	其中 \(K\) 为任一典范除子。证明 \(k\) 与典范除子的选择无关，且 \(W_{g-1}\) 在 \(u\mapsto k-u\) 下不变。证明存在 \(h\in J(X)\) 使 \(2h=k\)，并证明 \(W_{g-1}-h\) 关于原点对称。此题只用 Riemann--Roch 和 Abel 定理。
\end{exercise}
\begin{solution}
	两个典范除子之差为主除子，Abel 定理说明其积分像为零，所以 \(k\) 与选择无关。若 \(u\in W_{g-1}\)，取有效 \(D\)、\(\deg D=g-1\)，使 \(u=\operatorname{AJ}(D-(g-1)[q])\)。Riemann--Roch 此时给出
	\[
		h^0(\mathcal O(D))=h^0(\mathcal O(K-D)).
	\]
	左边至少为 \(1\)，故 \(K-D\) 线性等价于某个有效除子 \(D'\)。其次数仍为 \(g-1\)，并有
	\[
		A_q^{(g-1)}(D')
		=\operatorname{AJ}(K-D-(g-1)[q])=k-u.
	\]
	因此 \(k-W_{g-1}\subseteq W_{g-1}\)。映射 \(u\mapsto k-u\) 的平方为恒等，故包含实际为相等。

	在 \(J(X)=\mathbb C^g/\Lambda\) 中任取 \(k\) 的代表 \(\widetilde k\)，令 \(h=[\widetilde k/2]\)，就有 \(2h=k\)。若 \(v=u-h\in W_{g-1}-h\)，则
	\[
		-v=(k-u)-h\in W_{g-1}-h.
	\]
	所以这个平移关于原点对称。半点 \(h\) 未必唯一，差别是二扭点；对每种选择，所证对称性都成立。
\end{solution}

\begin{exercise}
	\label{g4:ex:0404:diagonal-period}
	设 \(\tau=\operatorname{diag}(\tau_1,\tau_2)\)，其中 \(\operatorname{Im}\tau_i>0\)。证明
	\[
		\theta((z_1,z_2),\tau)=\theta(z_1,\tau_1)\theta(z_2,\tau_2),
	\]
	并证明其零除子在某点有两个横截分支。利用\cref{g4:thm:0404:riemann-theta}和 Abel 嵌入，说明这样的 \(\tau\) 不可能是任何光滑紧连通亏格二曲面的归一化周期矩阵。本题使用正文注明的外引 theta 定理。
\end{exercise}
\begin{solution}
	对角矩阵使指数分成两个一维指数之和，绝对收敛允许交换双重求和，故得到乘积分解。对应环面也是
	\[
		\mathbb C^2/(\mathbb Z^2+\tau\mathbb Z^2)
		=E_{\tau_1}\times E_{\tau_2}.
	\]
	对每个一维环面，Riemann theta 定理说其 theta 除子是 \(W_0\) 的平移，并取约化结构；因此 \(\theta(z_i,\tau_i)\) 在环面上恰有一个零点 \(r_i\)，且为简单零点。在 \((r_1,r_2)\) 附近取坐标 \(u,v\)，两个因子分别为 \(u\) 与 \(v\) 乘以全纯单位，乘积的零方程因而等价于
	\[
		uv=0.
	\]
	它有两个不同的光滑分支，切线分别为两坐标轴，所以零除子不光滑。

	若这个矩阵来自光滑亏格二曲面 \(X\)，Riemann theta 定理又给出 \(\Theta=W_1+\kappa_q\)。但 \(W_1=A_q(X)\)，由\cref{g4:cor:0403:abel-embedding}是嵌入的光滑曲线，平移之后仍光滑，与刚得到的横截交点矛盾。因此对称且虚部正定虽足以构造带主极化的复环面，并不足以保证它是光滑亏格二曲面的归一化周期数据。
\end{solution}

\begin{exercise}
	\label{g4:ex:0405:origin-automorphisms}
	设 \(E:y^2=4x^3-Ax-B\) 光滑，\(O\) 为无穷远点。确定保持 \(O\) 的全纯自同构群，证明：当 \(j(E)\notin\{0,1728\}\) 时它为二阶循环群，当 \(j(E)=1728\) 时为四阶循环群，当 \(j(E)=0\) 时为六阶循环群。最后描述不要求保持 \(O\) 的全部全纯自同构。
\end{exercise}
\begin{solution}
	\cref{g4:thm:0405:j-classification}的系数比较已经说明，每个保持 \(O\) 的同构都写成
	\[
		(x,y)\longmapsto(u^2x,u^3y),\qquad u\ne0.
	\]
	在自同构情形，条件为
	\[
		u^4A=A,\qquad u^6B=B.
	\]
	若 \(AB\ne0\)，则 \(u^4=u^6=1\)，相除得 \(u^2=1\)，所以只有 \(u=\pm1\)。若 \(B=0\)，光滑性使 \(A\ne0\)，条件恰为 \(u^4=1\)。若 \(A=0\)，则 \(B\ne0\)，条件恰为 \(u^6=1\)。这些参数分别是二次、四次、六次单位根，乘法对应同构复合。参数不同确实给出不同同构：若 \(u^2=v^2\) 且 \(u^3=v^3\)，则 \(u=v\)。

	由 \(j=1728A^3/(A^3-27B^2)\)，判别式非零时 \(j=0\) 等价于 \(A=0\)，\(j=1728\) 等价于 \(B=0\)，所以恰得到所列三种情况。任意自同构 \(F\) 先减去点 \(F(O)\)，便成为保持 \(O\) 的自同构。因此全部自同构唯一写成
	\[
		P\longmapsto Q+\varphi(P),
		\qquad Q\in E,\quad \varphi\in\operatorname{Aut}(E,O).
	\]
	其中平移参数是 \(Q=F(O)\)，而有限群部分由上面的单位根给出。
\end{solution}

\begin{exercise}
	\label{g4:ex:0405:zeta-quasiperiods}
	设 \(\Lambda=\mathbb Z\lambda_1+\mathbb Z\lambda_2\)，且 \(\operatorname{Im}(\lambda_2/\lambda_1)>0\)。定义
	\[
		\zeta_\Lambda(z)=\frac1z+
		\sum_{\lambda\in\Lambda\setminus\{0\}}
		\left(\frac1{z-\lambda}+\frac1\lambda+\frac z{\lambda^2}\right).
	\]
	证明级数在 \(\mathbb C\setminus\Lambda\) 上正常收敛，\(\zeta_\Lambda'=-\wp_\Lambda\)，且存在常数 \(\eta_\lambda\)，使 \(\zeta_\Lambda(z+\lambda)-\zeta_\Lambda(z)=\eta_\lambda\)。证明 \(\eta_{\lambda+\mu}=\eta_\lambda+\eta_\mu\) 及
	\[
		\eta_{\lambda_1}\lambda_2-\eta_{\lambda_2}\lambda_1=2\pi\mathrm{i}.
	\]
	若 \([a]\ne[b]\in\mathbb C/\Lambda\)，用 \(\zeta_\Lambda\) 构造仅在这两点有单极点、留数分别为 \(1,-1\) 的第三类微分，并说明代表元的改变有何影响。
\end{exercise}
\begin{solution}
	在 \(|z|\le R\)、\(|\lambda|>2R\) 时，
	\[
		\frac1{z-\lambda}+\frac1\lambda+\frac z{\lambda^2}
		=-\frac{z^2}{\lambda^2(\lambda-z)},
	\]
	其绝对值至多为 \(2R^2|\lambda|^{-3}\)。格点计数已给出 \(\sum_{\lambda\ne0}|\lambda|^{-3}<\infty\)，所以级数在避开格点的紧集上正常收敛。每个格点是留数 \(1\) 的单极点，其余项全纯。逐项求导可由紧集上一致收敛与 Cauchy 公式保证，得到
	\[
		\zeta_\Lambda'(z)=-\frac1{z^2}
		+\sum_{\lambda\ne0}\left(-\frac1{(z-\lambda)^2}
		+\frac1{\lambda^2}\right)=-\wp_\Lambda(z).
	\]
	因此 \(\zeta_\Lambda(z+\lambda)-\zeta_\Lambda(z)\) 的导数为零，且同主部相消，差为常数 \(\eta_\lambda\)。连续平移两次便给出 \(\eta_{\lambda+\mu}=\eta_\lambda+\eta_\mu\)。

	取一个边界不经过格点的平移基本平行四边形，内部恰含一个格点。沿正向边界积分，留数定理给出 \(2\pi\mathrm{i}\)。将两条平行于 \(\lambda_1\) 的边配对，其贡献是 \(-\eta_{\lambda_2}\lambda_1\)；将两条平行于 \(\lambda_2\) 的边配对，其贡献是 \(\eta_{\lambda_1}\lambda_2\)。相加便得题中的关系，正向条件保证这里的符号为正。

	令
	\[
		\xi_{a,b}
		=\bigl(\zeta_\Lambda(z-a)-\zeta_\Lambda(z-b)\bigr)\,\mathrm{d}z.
	\]
	平移 \(z\) 一个格向量时，两项增加同一个准周期，差不变，所以微分下降到环面。单极点和留数由级数的局部主部给出，且没有其他极点。若把 \(a,b\) 分别换为 \(a+\lambda,b+\mu\)，微分增加 \((-\eta_\lambda+\eta_\mu)\,\mathrm{d}z\)，即一条全纯微分。再减去适当常数倍的 \(\mathrm{d}z\) 使指定 \(a\) 周期为零，就得到与代表元无关的归一化第三类微分，前提是归一化所用的实际环路也固定。
\end{solution}
\begin{exercise}
	\label{g4:ex:0405:singular-group-laws}
	考虑射影三次曲线的仿射方程
	\[
		C_{\mathrm n}:y^2=4(x-e)^2(x+2e)\quad(e\ne0),\qquad
		C_{\mathrm c}:y^2=4x^3,
	\]
	两者的无穷远点均记为 \(O\)。分别构造从 \(\mathbb P^1\) 到这两条曲线的全纯参数化，要求它在光滑部分一一对应，并确定奇点的全部原像。这些光滑参数曲线给出本题所指的归一化。证明两个光滑部分分别双全纯同构于 \(\mathbb C^\times\) 与 \(\mathbb C\)，且以 \(O\) 为单位元的弦切法分别对应乘法和加法。
\end{exercise}
\begin{solution}
	先看结点曲线。在 \(x\ne e\) 处令 \(t=y/(2(x-e))\)，方程给出 \(t^2=x+2e\)，所以
	\[
		x=t^2-2e,\qquad y=2t(t^2-3e).
	\]
	取 \(s\ne0\) 满足 \(s^2=3e\)。结点 \((e,0)\) 的原像恰为 \(t=s,-s\)，而 \(t=\infty\) 对应 \(O\)。若 \(t=u/v\)，参数化的齐次形式为
	\[
		[u:v]\longmapsto
		[(u^2-2ev^2)v:2u(u^2-3ev^2):v^3].
	\]
	这三个坐标从不同时为零，所以定义在整个 \(\mathbb P^1\) 上。有限光滑点的逆映射就是 \(t=y/(2(x-e))\)；在 \(O\) 附近，以 \(r=1/t\) 写成
	\[
		[\,r-2er^3:2-6er^2:r^3\,],
	\]
	其 \(X_1/Y\) 坐标导数在 \(r=0\) 为 \(1/2\)，故逆映射也在 \(O\) 全纯。因此光滑部分双全纯同构于 \(\mathbb P^1\setminus\{s,-s\}\)。再令
	\[
		q=\frac{t-s}{t+s},
	\]
	便把它识别为 \(\mathbb C^\times\)，并把 \(O\) 送到 \(q=1\)。两点 \(s,-s\) 在参数曲线上分离，正对应结点的两条分支。

	为核对群律，先取不经过结点且不竖直的直线 \(y=mx+b\)。其三个交点的参数 \(t_1,t_2,t_3\) 按重数计，是
	\[
		P(t)=2t^3-mt^2-6et+2me-b
	\]
	的三个根。因为 \(P(s)=P(-s)=-me-b\ne0\)，所以
	\[
		\prod_{j=1}^3\frac{t_j-s}{t_j+s}
		=\frac{P(s)}{P(-s)}=1.
	\]
	反射 \((x,y)\mapsto(x,-y)\) 对应 \(t\mapsto-t\)，从而对应 \(q\mapsto q^{-1}\)。因此把第三交点反射，所得参数正为前两点参数的乘积。重根已由上式按重数包含；竖直线的交点为 \(P,-P,O\)，也给出相同规则。通过两个光滑点的割线，或在光滑点的切线，不会额外通过结点：若直线通过结点，代入后的多项式已有 \(s,-s\) 两个根，最多只能再有一个光滑交点，且其重数为一。无穷远处的切线三次交于 \(O\)，与单位元相乘的规则相符。故整个光滑部分的弦切法就是 \(\mathbb C^\times\) 的乘法。

	尖点曲线的参数化为
	\[
		x=t^2,\qquad y=2t^3,\qquad
		[u:v]\longmapsto[u^2v:2u^3:v^3].
	\]
	尖点的唯一原像是 \(t=0\)，\(O\) 的原像为 \(t=\infty\)。在有限光滑点处，逆参数为 \(t=y/(2x)\)；在无穷远取 \(r=1/t\)，齐次表达为 \([r:2:r^3]\)，同样给出光滑的局部逆。因此光滑部分是 \(\mathbb P^1\setminus\{0\}\)，用
	\[
		w=\frac1t
	\]
	便得到到 \(\mathbb C\) 的双全纯映射，并有 \(w(O)=0\)。这里奇点只有一个参数原像，参数化在该点的导数为零，呈现尖点的一条分支。

	对不经过尖点的非竖直直线 \(y=mx+b\)，有 \(b\ne0\)，交点参数满足 \(2t^3-mt^2-b=0\)。一次项系数为零，所以按重数计，
	\[
		w_1+w_2+w_3
		=\frac{t_1t_2+t_1t_3+t_2t_3}{t_1t_2t_3}=0.
	\]
	反射对应 \(t\mapsto-t\)，即 \(w\mapsto-w\)，因此弦切法给出 \(w(P+Q)=w(P)+w(Q)\)。经过尖点的直线在 \(t=0\) 已有至少二重根，不能再包含两个光滑交点或在另一光滑点相切；竖直线与 \(O\) 的情形同前。于是加法规则覆盖整个光滑部分。两种群律均由这些明确坐标搬运，结合律和全纯性分别来自 \(\mathbb C^\times\) 与 \(\mathbb C\)。
\end{solution}